\documentclass[11pt,a4paper]{article}
\usepackage[margin=25mm]{geometry}
\usepackage{fontspec}
\setmainfont{DejaVu Serif}[ItalicFont=DejaVu Serif,ItalicFeatures={FakeSlant=0.17}]
\setmonofont{DejaVu Sans Mono}
\usepackage{amsmath,amssymb,amsthm,mathtools,booktabs,longtable,array,microtype,xcolor,xurl,enumitem}
\usepackage[colorlinks=true,linkcolor=blue!45!black,urlcolor=blue!45!black]{hyperref}
\hypersetup{pdftitle={Rank-nine fusion rings of multiplicity two},pdfauthor={FusionFrontier / Sebastien Palcoux}}
\newtheorem{theorem}{Theorem}[section]\newtheorem{proposition}[theorem]{Proposition}
\newtheorem{lemma}[theorem]{Lemma}\newtheorem{corollary}[theorem]{Corollary}
\theoremstyle{remark}\newtheorem{remark}[theorem]{Remark}
\newcommand{\FPdim}{\operatorname{FPdim}}\newcommand{\Rep}{\operatorname{Rep}}
\renewcommand{\Vec}{\operatorname{Vec}}\newcommand{\Irr}{\operatorname{Irr}}
\newcommand{\C}{\mathbb C}\newcommand{\Q}{\mathbb Q}\newcommand{\Z}{\mathbb Z}
\usepackage{fvextra}
\DefineVerbatimEnvironment{verbatim}{Verbatim}{breaklines,fontsize=\small}
\DefineVerbatimEnvironment{Highlighting}{Verbatim}{breaklines,commandchars=\\\{\}}
\providecommand{\tightlist}{\setlength{\itemsep}{0pt}\setlength{\parskip}{0pt}}
\setlength{\emergencystretch}{3em}
\title{Rank-nine fusion rings of multiplicity two\\\large Theoretical obstructions, coherent constructions and exact certification}
\author{FusionFrontier research prepared for S\'ebastien Palcoux\\\small AI-generated research; independent human verification and formalization pending}
\date{24 October 2026}
\begin{document}\maketitle
\begin{abstract}
We study the 2,083 supplied rank-nine based fusion rings of exact multiplicity
two. The current certified partition consists of 2,043 general complex
exclusions, 36 rings with unitary models and four unresolved rings. The proofs
combine formal-codegree and subring obstructions, signed sphericalization,
central induction, modular trace constraints, Morita theory and explicit
categorical constructions. A Galois restriction theorem simultaneously
excludes six rings by half-integral restriction multiplicities. A forced
central $\Rep(S_3)$ gives a further structural exclusion. Rank-one module
categories and the dimension-twenty Hopf classification settle three
additional rings, while explicit relative cocycles construct three integral
models. Every decision is bound to an original tensor, a proof dependency
and exact replay data. We classify existence of a categorification of a
based ring, without classifying all its associators.
\end{abstract}
\noindent\textbf{Verification scope.} These are AI-generated mathematical
claims. Separate mathematical audits and exact computational replays are
provided; independent human verification and Lean formalization are pending.
\tableofcontents

\section{The census, identifiers and accepted partition}
A fusion ring has a distinguished basis $1=X_0,X_1,\ldots,X_8$, nonnegative
integer constants $X_iX_j=\sum_kN_{ij}^kX_k$, an involution $i\mapsto i^*$,
and the usual unit and Frobenius reciprocity identities. Multiplicity means
$\max_{i,j,k}N_{ij}^k$. The source file through multiplicity two contains
2,225 tensors: 142 of multiplicity one and 2,083 of exact multiplicity two.
Its complete producer run finished all five duality strata. Enumeration
completeness is producer-reported; tensor validity and exact uniqueness are
separately checked. We do not repeat the entire enumeration here.

The label $R9\_2\_n$ denotes the $n$th \emph{exact-multiplicity-two} tensor
in source order, not necessarily physical line $n$ of the mixed source file.
Every ledger entry also supplies its physical line and SHA-256 of compact
JSON, without a trailing newline. The 142 multiplicity-one labels retain
explicit basis transports to the multiplicity-one classification; fourteen
of these transports are nonidentity. No literal unchanged-basis assertion is
needed. The original tensor, rather than a dimension profile, determines
case identity.

\begin{theorem}[Certified partition of the supplied exact-M2 census]
The accepted ledger partitions the 2,083 source rings into 2,043 rings
excluded from fusion categorification over $\C$, 36 admitting a unitary
fusion categorification, and four unresolved rings. Every exclusion has a
scope-valid theoretical argument and an exact certificate or an explicit
transfer from an excluded tensor. Every realization has categorical
existence data and a complete equality of its $9^3$ structure constants
with the source tensor under the recorded based permutation.
\end{theorem}
\begin{proof}
The finite partition is the exact union of the source-bound exclusion and
model records described below. The verifier checks all 2,083 tensor hashes,
the 568 explicit negative-subring transfers, the subsequent theoretical
certificates, and all 26,244 coefficients of the 36 positive tensors. It rejects any overlap between exclusion and
realization records. Each family-specific proof supplies the mathematical
implication used by its finite receipt. The model and exclusion arguments
are audited separately from these bookkeeping checks. Complete data and
replay commands appear in Section~\ref{sec:replay}.
\end{proof}

\paragraph{Subring-proof discipline.}
A closed based subring transfers an exclusion only when that subring has
an actual negative proof. Absence from a positive model catalog has no
logical force. In particular the rank-three half-$E_6$ ring
\[
 f^2=1,\qquad fx=xf=x,\qquad x^2=1+f+2x
\]
has a unitary categorification \cite{HH,ALW}. It is never used as a negative
subring. The audit withdraws 26 claims based on that omission, constructs
models for seven affected ambient rings, and replaces the proofs for
$R9\_2\_97,98,1489$ by their different $\mathrm{NG}(C_2^2,2)$ subrings.
For those subrings the established near-group restriction
$0<\kappa<|G|-1$ gives the exclusion. The corrected code contains an
explicit negative registry with tensor identities and proof dependencies.

\section{Uniform criteria and finite certificates}
\subsection{Formal codegrees, cyclotomicity and subrings}
For an irreducible representation $E$ of a based algebra let $e_E=\dim E$
and $f_E$ be its formal codegree. The fusion Casimir acts by $e_Ef_E$;
the degree factor must be retained in a noncommutative algebra. Formal
codegrees of a fusion category are $d$-numbers. If the minimal polynomial
of an algebraic integer $a$ is
$t^n+a_1t^{n-1}+\cdots+a_n$, the exact divisibilities
$a_n^i\mid a_i^n$ express the $d$-number test. Failure supplies a finite
integer witness. Character cyclotomicity is used with its proper scalar
character hypotheses; a nonabelian Galois-group witness cannot be replaced
by an approximate eigenvalue calculation \cite{ENO,Ostrik}.

For an integral ring a hypothetical categorification is pseudo-unitary.
Distinguished central dimensions $D/f_E$ restrict to sums of integral source
dimensions, hence are integers. In the noncommutative calculation the
central Casimir eigenvalue is divided by $e_E$ before imposing this
condition. Near-group and Tambara--Yamagami restrictions use their stated
group and coefficient hypotheses. Tensor closure, dual closure and a full
based permutation are checked before any proper-subring transfer.

The baseline dependency audit reviews all 588 original structural uses.
After removing the 26 invalid half-$E_6$ uses, the other 562 have direct
proofs or explicit lower negative bindings. The 454 rank-six through
rank-eight uses match excluded lower tensors exactly; 166 stored direct
witnesses replay; 24 small-rank uses are supplied with fresh direct
certificates. These counts describe overlapping audit classes, not a sum
of distinct excluded ambient rings.

\subsection{Central induction with representation degrees}
Let $\mathcal C$ be pseudo-unitary, $D=\FPdim(\mathcal C)$ and
$d_i=\FPdim(X_i)$. Write $F:\mathcal Z(\mathcal C)\to\mathcal C$ for
forgetting the half-braiding and $I$ for its right adjoint. The
distinguished simple $A_E\subset I(1)$ has
\[
 [I(1):A_E]=e_E,\qquad \dim A_E=D/f_E.
\]
For every central simple $Z$, its restriction column
$v_Z=([F(Z):X_i])_i$ is nonnegative integral, and
\[
 [I(X_i):Z]=v_{Z,i},\qquad \dim Z=\sum_iv_{Z,i}d_i.
\]
The full noncommutative induction identity is
\[
 F I(X)=\bigoplus_{Y\in\Irr(\mathcal C)}Y\otimes X\otimes Y^*.
\]
Consequently the columns satisfy the exact Gram identity
\[
 \sum_Zv_{Z,i}v_{Z,j}=[FI(X_i):X_j].
\]
In a commutative ring this is the usual Casimir Gram matrix; the simplification
is not used when the multiplication is noncommutative.

Every summand of $I(1)$ has twist one. The first and second twist traces in
$I(X_i)$, and the traces of cyclic rotations in tensor powers, give further
necessary equalities. Arithmetic and Galois conditions on central dimensions
and the integrality of modular character values restrict the finite column
domains. Retaining a larger domain only weakens an exclusion test.

\begin{proposition}[Exact separating certificate]
Suppose a finite rational system $Au=b$, $u\geq0$, is proved necessary for a
categorification. If an integer vector $y$ satisfies $A^ty\geq0$ and
$b^ty<0$, there is no such categorification.
\end{proposition}
\begin{proof}
A solution would give $b^ty=u^tA^ty\geq0$, a contradiction. Every entry of
the matrix and ray is checked exactly, after clearing denominators. A
failed numerical optimization or an unsuccessful search for $y$ has no
negative mathematical implication.
\end{proof}

\subsection{Signed reduction, neutral components and Morita theory}
A criterion formulated for pseudo-unitary categories is transferred to
arbitrary complex categories only through an established existence reduction.
Weak integrality supplies pseudo-unitarity directly. In the indicated
irrational cases, exact signed-sphericalization and character certificates
show that a hypothetical categorification has a pseudo-unitary Galois
conjugate of the same based ring. This is an equivalence of existence
questions, not an assertion that every original category is pseudo-unitary.
The signed audit checks the complete finite sign-character domains and
nonzero-product constraints. A signed gate alone never excludes a ring.

The graded and Morita arguments use actual invertible module categories.
For matrix conventions where source indices act on rows, module matrices
represent the opposite multiplication:
$M_iM_j=\sum_kN_{ji}^kM_k$. The linking-category Gram argument is applied
with this convention, including noncommutative cases. A feasible Gram
relaxation does not construct an invertible bimodule.

The accepted finite families, their whole-domain reductions and exact
rays are retained in the replay archive. The 36 intermediate exclusions
are grouped in Table~\ref{tab:families}. All their proofs and scope
reductions passed a separate audit. The folded $\mathrm{SO}(3)$ argument
for $R9\_2\_920$ includes the odd-boundary step: split a boundary strand
by the normalized trivalent vertex and retract by the bigon identity.
Thus spanning and relations extend from even boundaries without an
unsupported bending assertion.
\begin{small}\begin{longtable}{p{.36\textwidth}p{.56\textwidth}}
\caption{The 36 family exclusions. Listed integers are suffixes of $R9\_2\_n$.}\label{tab:families}\\
\toprule Family & Cases\\\midrule\endfirsthead\toprule Family & Cases\\\midrule\endhead
Signed nonzero products & 686,\allowbreak 690,\allowbreak 726,\allowbreak 728,\allowbreak 736,\allowbreak 1187,\allowbreak 1241,\allowbreak 1377,\allowbreak 1424\\
Joint integral center & 301,\allowbreak 352,\allowbreak 366,\allowbreak 385,\allowbreak 725,\allowbreak 727,\allowbreak 1028,\allowbreak 1134,\allowbreak 1238,\allowbreak 1448,\allowbreak 1465,\allowbreak 1501,\allowbreak 1712,\allowbreak 1778,\allowbreak 1989\\
Prime descent & 907\\
Solvable rank-two module & 1071,\allowbreak 1724\\
Joint quadratic center & 1569,\allowbreak 1886\\
Folded fermionic SO(3) & 920\\
Solvable central invertibles & 1356,\allowbreak 1933\\
Quadratic neutral block & 500,\allowbreak 1235\\
Fixed-block Morita obstruction & 396,\allowbreak 1265\\
\bottomrule\end{longtable}\end{small}


\section{A Galois restriction theorem}\label{sec:galois}
The following criterion reconstructs an entire central restriction from the
source algebra. It does not enumerate the rest of the center.
\begin{theorem}[Galois-unit restriction]
Let $R$ be a commutative fusion ring whose FP dimensions generate a real
quadratic field $K$. Suppose every formal codegree belongs to $K$, and
write $\sigma\ne1$ for its automorphism. Let $C_R$ be the regular Casimir
matrix. Its distinct irreducible minimal factors have degree one or two.
Define $p\in\Q[t]$ modulo their product by
\[
 p(t)\equiv1\pmod{t-a},\qquad
 p(t)\equiv t^2/n\pmod{t^2-at+n}.
\]
If $R$ has a pseudo-unitary categorification, the unit row of $p(C_R)$ is
a nonnegative integer vector with unit entry one.
\end{theorem}
\begin{proof}
Let $\chi_h$ run through all scalar irreducible characters, with formal
codegree $f_h$, and let $A_h\subset I(1)$ be their distinguished central
simples. Put $D=\FPdim(\mathcal C)$ and let $B$ correspond to
$\sigma\circ\FPdim$. Then
$\dim A_h=D/f_h$, $\dim B=D/\sigma(D)$, and $[F(B):1]=1$.
The central-character identity and symmetry of the unnormalized modular
$S$-matrix \cite{Deshpande} give
\begin{align*}
 (D/f_h)\chi_h(FB)&=S_{A_h,B}=S_{B,A_h}\\
 &=\frac{D}{\sigma(D)}(\sigma\circ\FPdim)(F A_h)
 =\frac{D}{\sigma(D)}\sigma(D/f_h)
 =\frac{D}{\sigma(f_h)}.
\end{align*}
Restriction multiplicities are integers, which justifies commuting
$\sigma$ with the dimension of $F A_h$. Hence
\[
 \chi_h(FB)=f_h/\sigma(f_h).
\]
For rational $f_h$ this is one; for a quadratic root with norm $n$ it is
$f_h^2/n$. Rational Chinese remaindering therefore constructs the displayed
$p$. All characters agree on $FB$ and $p(C_R)$, and the semisimple
commutative regular representation is faithful. Thus $p(C_R)$ is exactly
multiplication by $FB$, proving its required integrality and unit entry.
The characters themselves need not take their values in $K$; only their
formal codegrees do.
\end{proof}

\begin{corollary}
The rings $R9\_2\_114,128,1487,1635,1636,1959$ have no complex fusion
categorification.
\end{corollary}
\begin{proof}
Their FP field is $\Q(\sqrt3)$, with $D=96+48\sqrt3$. The forced
restrictions, in source order, are
\[
\begin{array}{c|c}
114&(1,0,0,0,\tfrac12,\tfrac12,1,0,0)\\
128&(1,0,0,\tfrac12,\tfrac12,\tfrac12,\tfrac12,0,0)\\
1487&(1,\tfrac12,\tfrac12,0,0,0,0,\tfrac12,\tfrac12)\\
1635&(1,\tfrac12,\tfrac12,0,0,0,0,0,1)\\
1636&(1,\tfrac12,\tfrac12,0,0,0,0,\tfrac12,\tfrac12)\\
1959&(1,\tfrac12,\tfrac12,\tfrac12,\tfrac12,0,0,0,0).
\end{array}
\]
Each contains nonintegral multiplicities. The exact independent replay
constructs the polynomial by a rational remainder linear system, checks
its values in every character component, verifies the full multiplication
matrix identity, and checks dimension $D/\sigma(D)=7+4\sqrt3$.
The separately audited signed gates reduce every general complex existence
question here to pseudo-unitary existence. The contradiction therefore has
general scope.
\end{proof}

\begin{remark}[Duality variant and noncommutative extension]
A dual-invariant scalar character labels a self-dual distinguished central
simple: $S_{A^*,Z}=S_{A,Z^*}$ and the character formula identify their
$S$-rows. In cases 1487 and 1959 a character of codegree 24 forces dimension
$4+2\sqrt3$. Its restriction must contain the unit and exactly one of a
pair of dual simples of dimension $2+2\sqrt3$, which is impossible.
For a noncommutative irreducible $E$ of degree $e_E$, the central eigenvalue
is $\chi_E/e_E$. The same proof gives $f_E/\sigma(f_E)$, while the Casimir
eigenvalue is $e_Ef_E$. The integer $e_E$ cancels in the Galois ratio.
The extension yields integral nonnegative vectors for 1848 and 1975 and
makes no exclusion for them. All eleven applicable commutative quadratic
positive controls pass the criterion.
\end{remark}

\section{A forced central symmetric category excludes 516}
\begin{proposition}
The ring $R9\_2\_516$ has no fusion categorification over $\C$.
\end{proposition}
\begin{proof}
Use its audited general-to-pseudo-unitary gate. The exact codegree data
force precisely one distinguished object of dimension one and one, $A$,
of dimension two in $I(1)$; every other such dimension exceeds four.
Its restriction is $F(A)=1+g$, with $g^2=1$. Uniqueness makes $A$ self-dual.
The unit multiplicity of $F(A)^2$ is two. Since all distinguished
multiplicities in $I(1)$ are one, $A^2$ has exactly one other distinguished
summand in addition to the unit. The dimension bound forces that summand
to be $A$, leaving an invertible $z$ with $F(z)=g$. Thus
\[
 A^2=1+z+A,\qquad zA=A,\qquad z^2=1.
\]
The verified character of $A$ has value one at $g$. The character formula
and balancing give $\theta_z=1$; all three objects have trivial double
braiding on their closed fusion subcategory. Positive dimensions identify
this symmetric category as Tannakian $\Rep(S_3)$.

Condense the central regular algebra $1+z$. Its image $1+g$ is connected,
so its module category $\mathcal E$ is a fusion category. The spherical
trace descends after division by two, giving its positive spherical
structure. Free and fixed source orbits under $g$ give simple dimensions
\[
 1^3,\quad d^3,\quad e^3,\qquad d=1+\sqrt3,\quad e=2+\sqrt3.
\]
The local $1+z$-modules in the central $\Rep(S_3)$ form $\Rep(C_3)$.
Under center descent \cite[Theorem 3.20]{DMNO}, its two nontrivial
invertibles forget separately to the unit of $\mathcal E$: their sum is
the free module on $F(A)=1+g$, namely two units. Distinct central structures
on the unit give a faithful $C_3$ of tensor automorphisms of the identity.
Their eigencharacters define a faithful dual $C_3$ grading, whose components
have equal FP dimension \cite{ENO}. Equating the rational and $\sqrt3$
coefficients, a component contains either three dimension-$d$ simples or
one each of dimensions $1,d,e$. The neutral component contains a unit,
so it has dimensions $1,d,e$.

Its rank-three fusion rules are forced by dimension and duality:
\[
 X^2=1+X+Y,\quad XY=X+2Y,\quad Y^2=1+2X+2Y.
\]
Its three distinguished restriction columns are $(1,0,0),(1,1,1),(1,2,2)$,
with dimensions $1,4+2\sqrt3,7+4\sqrt3$. Their twist-one contribution
in $I(X)$ is $M=18+10\sqrt3$. The total dimension of $I(X)$ is
$T=30+18\sqrt3$. Its first twist trace is zero, since $X\ne1$.
All other twists have modulus one, requiring $M\leq T-M$. But
$2M-T=6+2\sqrt3>0$, the contradiction.
An independently reconstructed labelled-$S$ ray against 14,128 columns
supplies an additional exact certificate, without being needed for this
structural proof.
\end{proof}

\section{The dimension-forty rank-one extension family}
\subsection{A fiber functor gives two exclusions}
\begin{proposition}
The rings $R9\_2\_1225$ and $R9\_2\_1673$ have no complex fusion
categorification.
\end{proposition}
\begin{proof}
In either ring the eight labels other than 5 form a fusion subring
$\mathcal D$ of dimensions $1^4,2^4$ and FP dimension 20. The other simple
$m$ satisfies
\[
 Xm=mX=d_Xm,\qquad m^2=\sum_Xd_XX.
\]
The additive category generated by $m$ is a rank-one left
$\mathcal D$-module category, equivalent to $\Vec$. Its action gives a
strong exact faithful tensor functor $\mathcal D\to\operatorname{End}(\Vec)
=\Vec$. Ordinary Tannaka reconstruction gives
$\mathcal D\simeq\Rep(H)$ for a semisimple Hopf algebra of dimension 20.
This uses neither a pivotal structure nor Hopf self-duality.

The complete dimension-twenty semisimple list consists of group algebras,
their duals, and the two exceptional Gelaki algebras $A_0,A_1$
\cite{ChengNg,Natale,AN}. Function Hopf algebras are pointed; group
representation rings are commutative. The two exceptional character rings
are also commutative, as the direct calculation below proves. Both source
neutral rings are noncommutative, witnessed respectively by the products
of labels $(1,2)$ and $(1,3)$. Thus none of the possibilities can occur.
\end{proof}

\subsection{Explicit Hopf data and a unitary model}
For an odd prime $p$ and $t=0,1$, write
$B(a,e,b)=x^a\delta_e g^b$, with $a\in\Z/p$ and $e,b\in\Z/2$.
The Andruskiewitsch--Natale construction \cite{AN} has
\begin{align*}
 B(a,e,b)B(c,f,d)&=\delta_{e,f}(-1)^{tebd}
 B(a+(-1)^bc,e,b+d),\\
 \Delta B(a,e,b)&=\sum_{s=0,1}B(a,s,b)\otimes B((-1)^sa,e-s,b),\\
 \epsilon B(a,e,b)&=\delta_{e,0}.
\end{align*}
Its four one-dimensional simples $L(e,s)$ have $x=1$ and
$g=(-1)^s i^{te}$. The remaining $p-1$ simples $V(e,j)$,
$1\leq j\leq(p-1)/2$, have dimension two and $x$-eigenvalues
$\zeta_p^{\pm j}$. Their characters are invariant under $a\mapsto-a$,
so inversion in the coproduct is invisible in every character convolution.
This proves commutativity of both exceptional character rings. Explicitly,
\begin{align*}
 L(e,s)L(f,r)&=L(e+f,s+r+tef),\\
 L(e,s)V(f,j)&=V(e+f,j),\\
 V(e,j)V(f,k)&=V(e+f,j+k)+V(e+f,j-k),
\end{align*}
where $V(e,-j)=V(e,j)$ and $V(e,0)=L(e,0)+L(e,1)$.

Put $\zeta=\zeta_p$. An explicit symmetric nondegenerate ordinary Hopf
pairing and its inverse copairing are
\begin{align*}
 P(B(a,e,b),B(c,f,d))&=\delta_{e,d}\delta_{f,b}\zeta^{ac}i^{tbd},\\
 \gamma&=\frac1p\sum_{a,c,e,b}\zeta^{-ac}i^{-teb}
 B(a,e,b)\otimes B(c,b,e).
\end{align*}
The positive Haar functional is
$h(B(a,e,b))=\delta_{a,0}\delta_{b,0}/2$.
Set $\lambda=\sqrt{4p}\,h$. Directly,
\[
 (\Delta\otimes1)\gamma=\gamma_{13}\gamma_{23},\qquad
 (1\otimes\Delta)\gamma=\gamma_{12}\gamma_{13},
\]
\[
 (S\otimes1)\gamma=(1\otimes S)\gamma,\qquad
 (\lambda\otimes\lambda)(1\otimes S)\gamma=1.
\]
The copairing is symmetric, $S^2=1$, and the distinguished grouplike is one.
These are precisely the conditions of \cite[Corollary 3.19]{DR}; rigidity
is supplied by their Section 3.6. Hence there is a coherent fusion category
$\Rep(A_t)+\Vec$ with the displayed odd rules. In their dual-Hopf
convention the self-duality map is $S\Gamma_0$, but their formula for the
associator inserts $S^{-1}$. Thus the actual associator copairing is the
$\gamma$ displayed above, without another antipode.

The compact involution is
\[
 B(a,e,b)^*=(-1)^{teb}B(-(-1)^ba,e,b).
\]
The Haar inner product makes these basis vectors orthogonal of squared
norm $1/2$. Each block of $\gamma$ is a Fourier bicharacter element times
a unitary and a phase. The odd Fourier associator is
\[
 \phi(B(a,e,b))=\frac{i^{-teb}}{\sqrt p}
 \sum_u\zeta^{-(-1)^eua}B(u,b,e),
\]
a permutation of normalized unitary Fourier blocks. The remaining
nonidentity associators are representations of $\gamma$, the Hopf
transport maps $h\otimes v\mapsto h_{(1)}\otimes h_{(2)}v$, and their
inverses. Haar invariance makes these transports unitary. In the odd-even-odd
case right multiplication composed with the antipode is a star
representation, with the corresponding opposite algebra understood.
All associators are therefore unitary.

For $p=5$, $t=1$ gives $R9\_2\_965$, with source-to-model map
$(0,2,3,6,7,8,5,4,1)$. The $t=0$ construction recovers $R9\_2\_21$ as a
positive control. The independent checker verifies both copairing identities,
normalization, all character convolutions, and all 729 coefficients.

\section{Three relative-cocycle models}
\subsection{A construction with its twisted character formula}
Let $G$ be a finite group, $H<G$, $\omega\in Z^3(G,U(1))$, and let
$\psi$ be normalized with $\delta\psi=\omega^{-1}|_H$ in the convention
\[
 \psi(a,b)\psi(ab,c)=\omega(a,b,c)\psi(b,c)\psi(a,bc).
\]
The subgroup algebra $A=\bigoplus_{h\in H}h$ in $\Vec_G^\omega$ is
connected and separable. Its bimodules form the unitary group-theoretical
category $\mathcal C(G,\omega,H,\psi)$ \cite{ENO,GHR}.
For a double-coset representative $u$ put $K_u=H\cap uHu^{-1}$.
On its degree-$u$ fiber, the stabilizer representation has multiplier
\[
 \alpha_u(a,c)=
 \frac{\psi(a,c)\omega(a,u,u^{-1}cu)}
 {\omega(a,c,u)\omega(u,u^{-1}au,u^{-1}cu)
 \psi(u^{-1}au,u^{-1}cu)}.
\]
This follows by writing the left and right module actions:
\[
 L_aL_c=\psi(a,c)\omega(a,c,v)^{-1}L_{ac},\quad
 R_cR_a=\omega(v,a,c)\psi(a,c)R_{ac},\quad
 R_cL_a=\omega(a,v,c)L_aR_c.
\]
One defines $\rho_u(a)=R_{u^{-1}au}^{-1}L_a$ and substitutes these identities.
For the relative tensor product the balancing relation is
\[
 R_h(m_u)\otimes n_v=\omega(u,h,v)m_u\otimes L_h(n_v).
\]
In a stabilizer trace at degree $uv$, a fixed coset contributes
\[
 \frac{\omega(u,h,v)}{\omega(a,u,v)\omega(u,v,b)}
 \chi_u^M(a)\chi_v^N(h),\quad
 h=u^{-1}au,\quad b=(uv)^{-1}a(uv).
\]
Projective character inner products recover fusion coefficients exactly.
The transport phase between homogeneous fibers is derived from the same
three module identities and recorded with a separate derivation in the
supplement. An untwisted convolution would be incorrect here.

\subsection{The model of dimension 120}
For $G=S_5$, $H=S_3\times C_2$, let $\varepsilon$ be parity and
$\alpha$ the supplied normalized $\mathbb F_2$ two-cocycle, defined by
a fixed section of the Schur double cover. Put
\[
 \omega(g,h,k)=(-1)^{\varepsilon(g)\alpha(h,k)}.
\]
The explicit unit-modulus cochain $\psi$ trivializes its restriction to
$H$. The independent replay reconstructs $S_5$ from permutations and
checks all $120^3$ two-cocycle equations and all $12^3$ algebra equations.
The three stabilizer orders are $12,2,4$, with projective irreducible
sizes $(1,1,1,1,2,2),(1,1),(2)$. The nondegenerate multiplier on the last
$V_4$ gives $M_2(\C)$. Simple dimensions are $1^4,2^2,6^3$, and the complete
tensor is $R9\_2\_92$ under $(0,3,1,5,7,6,4,2,8)$.

\subsection{Two models on the same group and subgroup}
Take $G=S_4\times C_2$, write $e$ for the external bit and $s$ for parity,
and set
\[
 x=((12),1),\quad y=((34),0),\quad H=\langle x,y\rangle\simeq V_4.
\]
Let $X,Y$ be its two coordinates and pull the explicit Schur cocycle
$\alpha$ back from $S_4$. The first associator and trivialization are
\[
 \omega_+=(-1)^{e\smile(\alpha+e^2)},\qquad
 \psi_+(a,b)=i^{X(a)Y(b)}.
\]
They construct $R9\_2\_392$ under $(0,1,2,8,7,3,5,6,4)$.

For the second model let $\Omega$ be the copied eighth-root-valued
$S_4$ cocycle from the lower-rank project. Its full table and original
source hash are supplied; all $24^4$ cocycle identities and every
transported table entry are independently checked. Set
\[
 \omega_- =\Omega\,(-1)^{e\smile\alpha+s\smile e^2}.
\]
In the order $(1,y,x,xy)$ of $H$, a trivialization
$\psi_-(a,b)=\zeta_8^{q(a,b)}$ is given by
\[
 (q(a,b))=\begin{pmatrix}0&0&0&0\\0&1&2&7\\0&4&0&0\\0&1&6&7\end{pmatrix}.
\]
The resulting tensor is $R9\_2\_1246$ under $(0,8,7,5,6,3,1,2,4)$.
For each construction the double-coset sizes are $4,4,16,16,8$, with
stabilizer orders $4,4,1,1,2$ and projective dimensions
$(1,1,1,1),(2),(1),(1),(1,1)$. This gives type $1^4,2^3,4^2$.
The independent checker verifies all $48^4$ cocycle equations, the
subgroup algebra and all tensor entries over the exact cyclotomic field.
Neither construction relies on an unproved classification of dimension 24.

\section{The half-\texorpdfstring{$E_6$}{E6} family and the model inventory}
\subsection{Products and associative zesting}
Let $\mathcal E$ be the unitary half-$E_6$ category and $f$ its specified
central fermion \cite{ALW}. The products with Ising, $\Rep(S_3)$,
$\mathrm{PSU}(2)_5$ and $\Vec(C_3)$ give respectively
$R9\_2\_71,79,499,2020$. Their tensors are direct Kronecker products.
For $R9\_2\_1218$ use the Ising parity on $\mathcal E\boxtimes\mathrm{Ising}$
and zest the $(1,1)$ product by $f$. The associator cochain
$\nu(a,b,c)=i^{abc}$ has coboundary $(-1)^{abcd}$ on parity bits, exactly
canceling the central fermion crossing. All sixteen coherence equations
are checked. The cochains and the initial categories are unitary.

\subsection{A central-fermion shear with two coherent actions}
On $\mathcal D=\mathcal E\boxtimes\Vec(C_2)$ define
$T(Z,h)=(f^hZ,h)$. Its tensorator moves the inserted $f^k$ left past $Z$
using the inverse half-braiding, and multiplies the adjacent powers of $f$.
The half-braiding identities prove its pentagon. Composition gives $T^2$
the identity tensorator multiplied by $(-1)^{hk}$, because $c_f=-1$.
For either $u=i$ or $u=-i$, the components $\eta_{(Z,h)}=u^h$ are a
monoidal isomorphism $T^2\to\mathrm{Id}$:
\[
 u^{h+k\bmod2}(-1)^{hk}=u^hu^k.
\]
The action equation $T\eta=\eta T$ holds because $T$ preserves $h$.
These are two coherent unitary $C_2$ actions.

The action fixes $(1,0),(f,0),(x,0),(x,1)$ and exchanges $(1,1),(f,1)$.
Its equivariantization has four invertibles, four dimension-$(1+\sqrt3)$
simples and one dimension-two object $Y$. Write $\epsilon$ for the acting
sign and choose lifts $a,b$ over $x,xh$. A normalized $t:fx\to x$ and
$r=tc_x^{-1}:xf\to x$ satisfy
\[
 r(t\otimes1_f)=-t(1_f\otimes r).
\]
This is naturality and the half-braiding hexagon with $c_f=-1$; it makes
$f$ fix the two lifts of $x$ and exchange the two lifts of $xh$.
The rules include
\[
 a^2=1+f+2a,\quad Y^2=1+f+\epsilon+f\epsilon,
\]
\[
 aY=b+\epsilon b,\quad bY=a+\epsilon a,\quad ab=Y+b+\epsilon b.
\]
The sign of $\eta$ changes the unit-pairing eigenvalue on the odd fixed
simple. Thus one choice has
$b^2=1+f\epsilon+a+\epsilon a$, and the other
$b^2=\epsilon+f+a+\epsilon a$. These give $R9\_2\_96$ and
$R9\_2\_1266$, including their distinct dualities. All coefficients,
action coherence and forgetful identities are independently reconstructed.

\subsection{All 36 positive tensors}
Table~\ref{tab:models} lists the accepted models. For group-theoretical
and equivariantization entries, the archive supplies the actual cocycles,
action data and projective stabilizer calculations, rather than an action
on the Grothendieck ring alone. The exceptional AH3 entry is matched to
its published categorical tensor. The WZW, subfactor, product and known
near-group entries use their specified unitary realizations.
The word ``model'' asserts existence for a based ring and makes no
claim of novelty up to tensor equivalence.
\begin{small}\setlength{\tabcolsep}{4pt}\begin{longtable}{p{.14\textwidth}p{.61\textwidth}p{.18\textwidth}}
\caption{All unitary models. The final column is the source-to-model permutation; indices are zero-based.}\label{tab:models}\\
\toprule Case & Construction & Basis map\\\midrule\endfirsthead\toprule Case & Construction & Basis map\\\midrule\endhead
R9\_2\_21 & Parity-twisted V4 equivariantization of TY(C5) & \texttt{013576428}\\
R9\_2\_71 & half-E6 x Ising & \texttt{043276518}\\
R9\_2\_79 & half-E6 x Rep(S3) & \texttt{043576218}\\
R9\_2\_92 & C(S5, sign cup Schur, S3 x C2, explicit trivialization) & \texttt{031576428}\\
R9\_2\_96 & C2 equivariantization of the central-fermion shear on half-E6 x Vec(C2) & \texttt{023567418}\\
R9\_2\_369 & Rep(PGL(2,7)) & \texttt{018245367}\\
R9\_2\_389 & Inversion orbifold of cyclic generalized Moore-Read category, dimension16 & \texttt{056241873}\\
R9\_2\_392 & C(S4 x C2, e cup (Schur + e\textasciicircum{}2), V4 graph subgroup, explicit trivialization) & \texttt{012873564}\\
R9\_2\_499 & half-E6 x PSU(2)\_5 & \texttt{035412678}\\
R9\_2\_515 & Diagonal-fermion quotient (Ising box PSU(2)\_10)\_{1+psi box V10} & \texttt{054321687}\\
R9\_2\_748 & TY(C3)\textasciicircum{}S3 by inversion, parity coherence 0 & \texttt{016724583}\\
R9\_2\_812 & Rep(V4 semidirect Dih(C9)) & \texttt{016725384}\\
R9\_2\_813 & Rep(V4 semidirect Dih(C3 x C3)) & \texttt{016724583}\\
R9\_2\_936 & Rep(PSL(2,8)) & \texttt{085167243}\\
R9\_2\_965 & Davydov--Runkel rank-one extension of Gelaki A1, dim(A1)=20 & \texttt{023678541}\\
R9\_2\_1016 & C(S4,1,V4nonnormal,1) & \texttt{067854123}\\
R9\_2\_1211 & TY(C3)\textasciicircum{}S3 by inversion, parity coherence 1 & \texttt{067245813}\\
R9\_2\_1218 & Fermionic associative C2 zest of half-E6 x Ising & \texttt{025768143}\\
R9\_2\_1246 & C(S4 x C2, exotic S4 associator times e cup Schur plus sign cup e\textasciicircum{}2, V4 graph subgroup, explicit trivialization) & \texttt{087563124}\\
R9\_2\_1266 & C2 equivariantization of the central-fermion shear on half-E6 x Vec(C2) & \texttt{067832154}\\
R9\_2\_1372 & Asaeda-Haagerup category AH3 & \texttt{034128576}\\
R9\_2\_1373 & (Fib box Fib)\textasciicircum{}S3, sign quotient acts by swapping factors & \texttt{045138762}\\
R9\_2\_1374 & D16 even category & \texttt{087246135}\\
R9\_2\_1396 & Pulled-back LPR-24 action along Dih(C9)->S3 & \texttt{067253814}\\
R9\_2\_1397 & Pulled-back LPR-24 action along Dih(C3 x C3)->S3 & \texttt{067245813}\\
R9\_2\_1557 & C(GL(3,2),1,S4,1) & \texttt{067124385}\\
R9\_2\_1857 & Spin(8)\_2 WZW category equivariantized by order-three Dynkin triality & \texttt{087213564}\\
R9\_2\_1979 & Izumi near-group (Q8,2) & \texttt{013467582}\\
R9\_2\_1983 & Ordinary-swap equivariantization of the equal half-E6 diagonal-fermion product & \texttt{013587642}\\
R9\_2\_2020 & half-E6 x Vec(C3) & \texttt{045871236}\\
R9\_2\_2029 & Triality equivariantization of the Tannakian-zested Spin(8)\_2 category & \texttt{087352164}\\
R9\_2\_2083 & Rep(C19 semidirect C3, u=7) & \texttt{012378564}\\
\bottomrule\end{longtable}\end{small}


{\setlength{\tabcolsep}{3pt}
\hypertarget{cyclic-internal-end-algebras-and-multiplicity}{%
\section{Cyclic internal-end algebras and
multiplicity}\label{cyclic-internal-end-algebras-and-multiplicity}}

Date: 24 October 2026. Scope: arbitrary fusion categories over C. No
pivotal, spherical, pseudo-unitary, unitary, or braided hypothesis
occurs.

\hypertarget{theorem}{%
\subsection{Theorem}\label{theorem}}

Let D be a fusion category, M a finite semisimple indecomposable left
D-module category, and x a simple object of M. Suppose its internal end
algebra has underlying object

\begin{verbatim}
A = [x,x]\_D = direct\_sum\_{g in H} g,
\end{verbatim}

where H is a finite cyclic group of invertible objects, each with
multiplicity one. Then y tensor x is multiplicity-free for every simple
object y of D.

\hypertarget{proof}{%
\subsubsection{Proof}\label{proof}}

For each g in H, the nonzero morphism g tensor x -\textgreater{} x is an
isomorphism. Consequently every component g tensor h -\textgreater{} gh
of the multiplication of A is nonzero: it is induced by composition of
these isomorphisms. This is not an assumption about an arbitrarily
chosen algebra structure.

The internal-Hom equivalence of Ostrik, Theorem 1, identifies M with the
category of right A-modules in D, sends x to A, and sends y tensor x to
the free module y tensor A. Therefore their endomorphism algebras agree.
By free-module adjunction,

\begin{verbatim}
End\_A(y tensor A) = Hom\_D(y, y tensor A)
                  = direct\_sum\_{g in K} Hom\_D(y,y tensor g),
\end{verbatim}

as vector spaces, where K = \{g in H : y tensor g is isomorphic to y\}.
K is a subgroup of H, hence cyclic. Every displayed summand is
one-dimensional. Choose nonzero isomorphisms u\_g:y -\textgreater{} y
tensor g. The associated free-module endomorphism is the composite

\begin{verbatim}
y tensor A --u\_g tensor id--> (y tensor g) tensor A
           --associator--> y tensor (g tensor A)
           --id tensor multiplication--> y tensor A.
\end{verbatim}

Call this endomorphism e\_g. The multiplication of A and the u\_g are
isomorphisms on the relevant simple summands, so composition satisfies
e\_g e\_h = c(g,h) e\_\{gh\} (or the reversed product according to
convention), with c(g,h) nonzero. In particular each e\_g is invertible.
If t generates K of order n, the powers of e\_t span all these
homogeneous summands and e\_t\^{}n = c id for c nonzero. Thus

\begin{verbatim}
End\_A(y tensor A) is isomorphic to C[T]/(T\^n-c).
\end{verbatim}

This algebra is commutative. In a semisimple C-linear category the
endomorphism algebra of direct\_sum\_i m\_i z\_i, for distinct simples
z\_i, is the product of Mat\_\{m\_i\}(C); it is commutative precisely
when every m\_i is at most one. This proves the theorem. The case
K=\{1\} is included.

The proof does not invoke H\^{}2(H,C*)=0, a fiber functor on D, a
classification of Tambara-Yamagami categories, or a normalization of
their associators. The cyclic hypothesis is essential: for V4 a twisted
group algebra can be Mat\_2(C).

\hypertarget{four-rank-nine-applications}{%
\subsection{Four rank-nine
applications}\label{four-rank-nine-applications}}

All labels are zero-based source-tensor labels. In every row, the
invertible group is G=\{0,1,2,8\}, generated by label 1 with powers
(0,1,8,2). In particular it is C4, not V4.

\begin{longtable}[]{@{}lllll@{}}
\toprule\noalign{}
Ring & D simples & y & x & x' \\
\midrule\noalign{}
\endhead
\bottomrule\noalign{}
\endlastfoot
R9\_2\_1025 & 0,1,2,6,8 & 6 & 4 & 5 \\
R9\_2\_1709 & 0,1,2,7,8 & 7 & 5 & 6 \\
R9\_2\_1759 & 0,1,2,6,8 & 6 & 3 & 4 \\
R9\_2\_1988 & 0,1,2,7,8 & 7 & 5 & 6 \\
\end{longtable}

For each tensor the listed D is a fusion subring with rules g y = y g =
y and y\^{}2 = sum\_\{g in G\} g. Its hypothetical categorical
realization is a fusion subcategory of any hypothetical categorification
of the full ring. The additive category M spanned by x,x' is a D-module
subcategory: each g fixes both objects and

\begin{verbatim}
y tensor x = 2 x',     y tensor x' = 2 x.
\end{verbatim}

It is indecomposable since y connects its two simples. The internal end
of x in D is exactly sum\_G g: the four invertibles each fix x, and y
has no diagonal action on x. The theorem forces y tensor x to be
multiplicity-free, contrary to 2 x'. Equivalently its endomorphism
algebra would have to be both commutative and Mat\_2(C).

Conclusion: all four rings have no categorification over C.

\hypertarget{exact-certificate-and-provenance}{%
\subsection{Exact certificate and
provenance}\label{exact-certificate-and-provenance}}

\texttt{cyclic\allowbreak{}\_\allowbreak{}tensor\allowbreak{}s.\allowbreak{}json}
contains the four full 9-by-9-by-9 tensors copied from the accepted
phase-three corpus, without alteration.
\texttt{verify\allowbreak{}\_\allowbreak{}cyclic\allowbreak{}\_\allowbreak{}obstru\allowbreak{}ction.\allowbreak{}py}
checks their SHA-256 values, all fusion ring axioms, C4 multiplication,
dual closure and multiplication closure of D, the complete action of D
on M, the internal-end support, and the off-diagonal multiplicity two.
The output is
\texttt{cyclic\allowbreak{}\_\allowbreak{}obstru\allowbreak{}ction\_\allowbreak{}certif\allowbreak{}icate.\allowbreak{}json}.
These finite checks verify the hypotheses of the theorem; they do not
replace the categorical proof.

\hypertarget{primary-source}{%
\subsection{Primary source}\label{primary-source}}

V. Ostrik, *Module categories, weak Hopf algebras and modular
invariants*, Transform. Groups 8 (2003), 177--206, arXiv:math/0111139,
Section 3.2, Lemma 5 and Section 3.3, Theorem 1. The proof explicitly
uses the internal end of a chosen object and the free-module adjunction,
and imposes no unitarity assumption.

\url{https://arxiv.org/pdf/math/0111139}

The cyclic obstruction above is a direct algebraic consequence proved
here; no claim of priority or literature novelty is made.

\hypertarget{a-nonsolvable-dimension-sixty-neutral-component}{%
\section{A nonsolvable dimension-sixty neutral
component}\label{a-nonsolvable-dimension-sixty-neutral-component}}

Date: 24 October 2026. General-complex scope, with no braiding or
unitarity assumptions on the hypothetical categorification.

\hypertarget{theorem-1}{%
\subsection{Theorem}\label{theorem-1}}

Let D be a fusion category with simple FP dimensions 1,1,2,2,5,5.
Assume:

\begin{enumerate}
\def\labelenumi{\arabic{enumi}.}
\tightlist
\item
  D equals its adjoint subcategory (its universal grading is trivial).
\item
  The nontrivial invertible g fixes the two dimension-two simples and
  exchanges the two dimension-five simples.
\end{enumerate}

Then every faithful finite-group extension of D is group-theoretical,
and in particular is integral.

\hypertarget{step-1-d-cannot-be-solvable}{%
\subsubsection{Step 1: D cannot be
solvable}\label{step-1-d-cannot-be-solvable}}

Suppose D were solvable. By ENO Proposition 4.4 it admits a construction
from Vec by successive extensions and equivariantizations by cyclic
groups of prime order. The last step cannot be an extension, since D has
no nontrivial faithful grading. Thus D = E\^{}\{C\_p\} for a fusion
category E, and the canonical central Rep(C\_p) embeds into D under
forgetting its central structure. D has exactly two invertibles, so p=2,
and the nontrivial character has underlying object g.

De-equivariantize by this central Rep(C2), or equivalently take modules
over its regular central algebra A=1+g. This is a fusion category of
dimension 30. Its simple dimensions can be determined without knowing
any associator of D. The free orbit \{1,g\} gives one invertible; the
free orbit consisting of the two dimension-five objects gives one simple
of dimension five. Each fixed dimension-two object gives two distinct
simple A-modules of dimension one. Indeed, the endomorphism algebra of
its free module has dimension two and is semisimple, hence is C+C; the
two simple summands each have underlying object equal to the original
dimension-two object. FP dimension after passage to A-modules is the
underlying FP dimension divided by FPdim(A)=2. These six objects exhaust
the simples since every simple A-module is a summand of a free module.
Therefore the de-equivariantization has dimensions 1\^{}5,5.

Its invertibles form the cyclic group C5. Writing z for its unique
noninvertible simple, rigidity and dimensions force

\begin{verbatim}
hz=z=zh,    z\^2 = sum\_{h in C5} h + 4z.
\end{verbatim}

But Etingof--Gelaki--Ostrik, Corollary 7.4, applies to arbitrary complex
fusion categories with n cyclic invertibles and one noninvertible of
dimension n.~It forces n+1 to be a prime power. Here n=5 and 6 is not a
prime power. This contradiction proves that D is not solvable.

\hypertarget{step-2-d-is-morita-equivalent-to-vec_a5omega}{%
\subsubsection{Step 2: D is Morita equivalent to
Vec\_\{A5\}\^{}omega}\label{step-2-d-is-morita-equivalent-to-vec_a5omega}}

ENO Theorem 9.16 says that every fusion category of dimension 60 is
weakly group-theoretical. Hence D is categorically Morita equivalent to
a nilpotent fusion category N of the same dimension. Choose a tower of
faithful group extensions constructing N from Vec. The product of the
orders of the grading groups is 60.

If all grading groups are solvable, their normal series refine this
tower into cyclic-prime extensions; N and hence D are solvable, contrary
to Step 1. Thus at least one grading group is nonsolvable. Every proper
divisor of 60 is the order of only solvable finite groups; the unique
nonsolvable group of order 60 is A5. These standard finite-group facts
can also be expressed using the smallest nonabelian simple group A5, of
order 60, and a composition series. The nonsolvable grading group
therefore has order 60, and every other extension step is trivial.
Consequently N is a pointed category Vec\_\{A5\}\^{}omega for some
3-cocycle omega. No value or trivialization of omega is assumed.

\hypertarget{step-3-all-extensions-of-vec_a5omega-are-pointed}{%
\subsubsection{Step 3: all extensions of Vec\_\{A5\}\^{}omega are
pointed}\label{step-3-all-extensions-of-vec_a5omega-are-pointed}}

Let P=Vec\_\{A5\}\^{}omega, and let B be an invertible P-bimodule
category. As a left P-module category it has pointed dual (indeed the
right action identifies its dual with P, up to the harmless opposite
convention). By the classification of pointed-dual module categories,
Naidu--Nikshych Example 2.2, such a module category is described by a
normal abelian subgroup H of A5 and compatible cochain data. Since A5
has no nontrivial normal abelian subgroup, H=\{1\}. Therefore B is the
regular left P-module category and has 60 simple objects, transitively
permuted by the 60 invertibles of P.

Each homogeneous component of a faithful group extension of P is an
invertible P-bimodule category. Consequently each component has these 60
simple objects, all with the same FP dimension d, because the
invertibles of P permute them transitively. A faithful homogeneous
component has total FP dimension FPdim(P)=60. Hence 60*d\^{}2=60 and
d=1. Every simple object of the extension is invertible: the extension
is pointed. This argument applies to every omega and every coherent
extension, not merely to untwisted or preselected models.

\hypertarget{step-4-transport-the-extension-through-morita-equivalence}{%
\subsubsection{Step 4: transport the extension through Morita
equivalence}\label{step-4-transport-the-extension-through-morita-equivalence}}

If C is any faithful finite-group extension of D, ENO Lemma 3.4
constructs an extension of P which is Morita equivalent to C. It does
this by replacing D with its category of bimodules over an
indecomposable algebra and applying the same construction to C; the
grading survives because the algebra lies in the neutral component. Step
3 makes the resulting category pointed. Thus C is group-theoretical.
Group-theoretical fusion categories are integral: the standard
double-coset description of simples gives dimensions

\begin{verbatim}
[H : H intersect tHt\^{-1}] \* dim(projective irreducible representation),
\end{verbatim}

which are integers. This proves the theorem.

\hypertarget{application-to-two-remaining-rank-nine-rings}{%
\subsection{Application to two remaining rank-nine
rings}\label{application-to-two-remaining-rank-nine-rings}}

Both tensors have global FP dimension 120. Their integral objects form a
dimension-60 fusion subring D and the other three objects each have
dimension 2*sqrt(5). The tensor verifies a faithful C2 grading with D as
neutral component. In the notation (1,g,a,b,u,v) of dimensions
(1,1,2,2,5,5), the neutral labels are:

\begin{longtable}[]{@{}lll@{}}
\toprule\noalign{}
Ring & Neutral labels (1,g,a,b,u,v) & Odd labels \\
\midrule\noalign{}
\endhead
\bottomrule\noalign{}
\endlastfoot
R9\_2\_585 & (0,1,7,8,3,4) & 2,5,6 \\
R9\_2\_1303 & (0,4,5,6,7,8) & 1,2,3 \\
\end{longtable}

The neutral rules satisfy the theorem's two conditions. In particular
the adjoint closure is all six labels, while g fixes a,b and exchanges
u,v. A hypothetical categorification would therefore be an integral
extension by the theorem. Its three dimensions 2*sqrt(5) contradict
integrality.

Conclusion: R9\_2\_585 and R9\_2\_1303 have no complex categorification.

This proof does not assume that a fusion-ring match to a known
dimension-60 model classifies its categorifications. Step 2 establishes
the necessary Morita form for every hypothetical categorification.

\hypertarget{exact-certificate}{%
\subsection{Exact certificate}\label{exact-certificate}}

\texttt{d120\_\allowbreak{}tensor\allowbreak{}s.\allowbreak{}json}
contains the two unchanged full source tensors.
\texttt{verify\allowbreak{}\_\allowbreak{}dimens\allowbreak{}ion120\allowbreak{}\_\allowbreak{}obstru\allowbreak{}ction.\allowbreak{}py}
checks all fusion-ring axioms, all FP dimension equations exactly in
Q(sqrt(5)), the neutral subcategory, faithful grading, its adjoint
closure, and the complete permutation of simples by g. Its output is
\texttt{dimens\allowbreak{}ion 12\allowbreak{}0\_\allowbreak{}obstru\allowbreak{}ction\_\allowbreak{}certif\allowbreak{}icate.\allowbreak{}json}.

\hypertarget{primary-references-and-theorem-scopes}{%
\subsection{Primary references and theorem
scopes}\label{primary-references-and-theorem-scopes}}

* P. Etingof, D. Nikshych and V. Ostrik, *Weakly group-theoretical and
solvable fusion categories*, Adv. Math. 226 (2011), 176--205,
arXiv:0809.3031v2. Definition 1.1; Proposition 4.4; Proposition 4.5;
Theorem 9.16 (dimension 60); Lemma 3.4 (Morita transport of extensions).
\url{https://arxiv.org/pdf/0809.3031} * P. Etingof, S. Gelaki and V.
Ostrik, *Classification of fusion categories of dimension pq*, IMRN
2004, no. 57, 3041--3056, arXiv:math/0304194v2, Corollary 7.4. The
statement concerns all complex fusion categories and is not a
unitary-only classification. \url{https://arxiv.org/pdf/math/0304194} *
D. Naidu and D. Nikshych, *Lagrangian subcategories and braided tensor
equivalences of twisted quantum doubles of finite groups*, Commun. Math.
Phys. 279 (2008), 845--872, arXiv:0705.0665v2, Example 2.2 and Section
4.2. The normal-abelian-subgroup condition covers arbitrary 3-cocycles
omega. \url{https://arxiv.org/pdf/0705.0665}

The theorem is proved here by combining the stated results. No claim of
priority or exhaustive literature novelty is made.

\hypertarget{cyclic-generalized-tambarayamagami-constructions}{%
\section{Cyclic generalized Tambara--Yamagami
constructions}\label{cyclic-generalized-tambarayamagami-constructions}}

\hypertarget{result}{%
\subsection{Result}\label{result}}

The rings R9\_2\_746 and R9\_2\_1209 have explicit unitary
categorifications. They are the two parity-coherence inversion
equivariantizations of a 12-simple-object generalized Tambara--Yamagami
category with invertible group C9 and three simple objects of dimension
sqrt(3). This category is constructed below by an explicit central C3
condensation, rather than inferred from vanishing of a group cohomology
obstruction.

All half-braidings in this note use the orientation
\texttt{c\_\allowbreak{}z(X): \allowbreak{}z tens\allowbreak{}or X -\allowbreak{}> X te\allowbreak{}nsor z\allowbreak{}}.
This is inverse to the orientation used for central structures in GNN
Section 2.2; keeping this convention explicit avoids a sign reversal in
the bicharacter.

\hypertarget{starting-unitary-categories-and-central-object}{%
\subsection{Starting unitary categories and central
object}\label{starting-unitary-categories-and-central-object}}

Let omega=exp(2 pi i/3), zeta=exp(2 pi i/9). Take the unitary category
T=TY(C3,chi,1/sqrt(3)), chi(a,b)=omega\^{}(ab), with the skeletal
associators of GNN Section 4.1. Its objects are a in C3 and m, with
m\^{}2=sum\_a a. The central lift z of the invertible a=1 has right
half-braiding

\begin{verbatim}
c\_z(b)=omega\^b,   c\_z(m)=omega.
\end{verbatim}

The equations are
\texttt{c\_\allowbreak{}z(m)\textasciicircum{}2\allowbreak{}=chi(1\allowbreak{},1)\textasciicircum{}(-\allowbreak{}1)}
and
\texttt{c\_\allowbreak{}(z ten\allowbreak{}sor z'\allowbreak{})(m)=c\allowbreak{}\_\allowbreak{}z(m)c\_\allowbreak{}z'(m)c\allowbreak{}hi(a,a\allowbreak{}')\textasciicircum{}(-\allowbreak{}1)}.
Thus z has order3 as a central object, and c\_z(z)=omega. These formulas
follow by substitution in the displayed TY associators; they are also
the inverse-half-braiding form of the central invertibles in GNN
Proposition 4.1. No braiding on T itself is asserted.

Let P=Vec(C9,Omega), with representatives0,\ldots,8 and

\begin{verbatim}
Omega(a,b,c)=zeta\^(3a floor((b+c)/9)).
\end{verbatim}

This is a normalized unitary 3-cocycle. The underlying invertible h=3
has right half-braiding beta(x)=zeta\^{}(2x). Its central square has
underlying object6 and half-braiding zeta\^{}(4x). Its central cube has
half-braiding

\begin{verbatim}
zeta\^(6x) Omega(x,6,3)=zeta\^(9x)=1.
\end{verbatim}

In the product category B=T boxtimes P, the central object

\begin{verbatim}
g=(z,3)
\end{verbatim}

has order3. Its self-braiding is omega*zeta\^{}6=1. The associator on
its three powers is identically1, since Omega(3a,3b,3c)=1. Their
braiding is also identically1, because the self-braiding of a generator
is1 and this is a pointed category with trivial associator. They
therefore give an explicit central Tannakian Rep(C3). Set A=1+g+g\^{}2,
with the ordinary group-algebra multiplication (all nonzero
multiplication scalars1). It is a commutative separable central algebra
with the usual unitary normalization. The right A-module category D=B\_A
is a unitary fusion category of dimension (6*9)/3=18.

The action of g on every simple of B is free: its C9 charge shifts by3.
Hence every free simple A-module is simple, and these are indexed by the
12 orbits. Writing the nine invertibles as u\^{}a, a in C9, and the
three other simples as X\_r, r in C3, gives

\begin{verbatim}
u\^a u\^b = u\^(a+b),
u\^a X\_r = X\_r u\^a = X\_(r+a),
X\_r X\_s = sum\_(h=0,1,2) u\^(r+s-3h).
\end{verbatim}

For example, use orbit representatives (0,a) and (m,r); the relation
(1,3)=1 in the quotient identifies (h,r+s) with (0,r+s-3h). The
dimensions are1\^{}9,(sqrt3)\^{}3.

\hypertarget{a-coherent-involution-preserving-the-algebra}{%
\subsection{A coherent involution preserving the
algebra}\label{a-coherent-involution-preserving-the-algebra}}

On T use the canonical orthogonal inversion autoequivalence
F(a)=-a,F(m)=m. The tensorator is the canonical identity map, including
the permutation of summands of m\^{}2. This is an actual involutive
tensor autoequivalence, as in GNN Proposition 4.8.

On P use F(x)=-x with tensorator

\begin{verbatim}
J(a,b)=zeta\^mu(a,b),  mu(a,b)=-3a\*1\_(b>0) modulo 9.
\end{verbatim}

Writing w for the exponent of Omega, direct substitution gives

\begin{verbatim}
mu(a,b)+mu(a+b,c)-mu(b,c)-mu(a,b+c)
  = w(-a,-b,-c)-w(a,b,c) modulo 9,
mu(a,b)+mu(-a,-b)=0 modulo 9.
\end{verbatim}

Consequently F\^{}2=Id as a tensor functor with identity square
coherence. On T it sends the central z to z\^{}2. On P the transported
half-braiding of h=3 at the new argument x is

\begin{verbatim}
beta(-x) J(3,-x)/J(-x,3) = zeta\^(4x),
\end{verbatim}

the half-braiding of the inverse central object h\^{}2. Thus F(g)=g\^{}2
as central objects. On powers of g the tensorator is identically1
because mu(3a,3b)=0. The permutation g\^{}i -\textgreater{} g\^{}(-i)
therefore defines an actual algebra isomorphism F(A)-\textgreater A
whose square is identity, respecting both multiplication and the central
structure. Hence F descends to a coherent unitary C2 action on D. It
acts by

\begin{verbatim}
u\^a -> u\^(-a),    X\_r -> X\_(-r).
\end{verbatim}

There is a second coherent action with the same tensor functor: multiply
the square coherence by the parity natural automorphism, +1 on pointed
simples and -1 on X\_r. Parity is monoidal, invariant under F, and
trivial on A; hence it obeys the C2 action coherence before and after
descent. Both actions consist entirely of unitary structure maps.

\hypertarget{exact-equivariant-fusion-tensors}{%
\subsection{Exact equivariant fusion
tensors}\label{exact-equivariant-fusion-tensors}}

Let E\_0 and E\_1 be the equivariantizations for the identity and
negative odd square coherence, respectively. They are unitary categories
of dimension 36. Inversion fixes the unit and X\_0, and has four
two-element pointed orbits and one two-element nonpointed orbit. Thus
each E has simple dimensions 1\^{}2,2\^{}4,(sqrt3)\^{}2,2sqrt3 and rank
9.

The two fixed objects give two equivariant lifts each. Every free orbit
gives one simple. The fusion multiplicity for a fixed output is
(dim+trace)/2; for a free output orbit it is the multiplicity at either
representative. Only fixed input objects contribute to the trace. The
only nontrivial fixed-input fusion line is X\_0 X\_0 -\textgreater1.
This line is inherited from (m,0)(m,0)-\textgreater(0,0), where the
tensorator acts identically. In E\_0 its trace is the product of the
equivariant signs. In E\_1 each fixed odd equivariant structure squares
to-1, so their product contributes the extra factor-1. This determines
every tensor coefficient without choosing unknown F-symbols on D.

The standalone generator
\texttt{cyclic\allowbreak{}\_\allowbreak{}ty\_\allowbreak{}models\allowbreak{}.\allowbreak{}py}
constructs both tensors from these orbit/trace formulas, verifies the
unit, duality, Frobenius reciprocity and all associativity equations,
and matches every one of the729 original census coefficients. The
source-to-model permutations are

\begin{longtable}[]{@{}lll@{}}
\toprule\noalign{}
Source & Coherence & Basis permutation \\
\midrule\noalign{}
\endhead
\bottomrule\noalign{}
\endlastfoot
R9\_2\_746 & identity & {[}0,1,6,7,2,5,3,8,4{]} \\
R9\_2\_1209 & negative odd square & {[}0,6,7,2,5,3,8,1,4{]} \\
\end{longtable}

\texttt{cyclic\allowbreak{}\_\allowbreak{}ty\_\allowbreak{}models\allowbreak{}.\allowbreak{}json}
contains both full tensors, the orbit labels, the intermediate rank 12
tensor, and the complete exact parameter checks. The finite checks
corroborate explicit categorical data; associativity of the fusion
tensor alone is not used as an existence argument.

\hypertarget{references}{%
\subsection{References}\label{references}}

\begin{itemize}
\tightlist
\item
  D. Tambara and S. Yamagami, Tensor categories with fusion rules of
  self-duality for finite abelian groups, J. Algebra209(1998),692--707.
\item
  S. Gelaki, D. Naidu and D. Nikshych, Centers of graded fusion
  categories, Algebra \& Number Theory3(2009),959--990; arXiv:0905.3117.
  Sections 2.4--2.5 give equivariantization/de-equivariantization;
  Section 4.1 gives the explicit TY associators; Proposition 4.1
  describes central invertibles and Proposition 4.8 gives the orthogonal
  action. Primary source checked: \url{https://arxiv.org/pdf/0905.3117}.
\end{itemize}

This settles existence only; no classification or novelty claim about
all tensor-equivalence classes with these fusion rings is made.

\hypertarget{uniform-family}{%
\subsection{Uniform family}\label{uniform-family}}

The same construction works with3 replaced by any odd n\textgreater1.
Put zeta=exp(2 pi i/n\^{}2), chi(a,b)=zeta\^{}(n ab) on Cn, choose
\texttt{c\_\allowbreak{}z(m)=c\allowbreak{}hi(1,1\allowbreak{})\textasciicircum{}((n-\allowbreak{}1)/\allowbreak{}2)}
(whose square is chi(1,1)\^{}(-1)), and use

\begin{verbatim}
Omega(a,b,c)=zeta\^(n a floor((b+c)/n\^2)),
h=n, beta(x)=zeta\^((n-1)x),
J(a,b)=zeta\^(-n a 1\_(b>0)).
\end{verbatim}

The n-th central power of h has half-braiding
\texttt{beta(x\allowbreak{})\textasciicircum{}n ze\allowbreak{}ta\textasciicircum{}(n \allowbreak{}x)=1}.
The diagonal generator's self-braiding is1. Both the associator and
inversion tensorator restrict trivially to its powers. Inversion sends
its central lift to its inverse because
\texttt{-\allowbreak{}(2n-\allowbreak{}1)=(n-\allowbreak{}1)\textasciicircum{}2 m\allowbreak{}odulo \allowbreak{}n\textasciicircum{}2}.
This gives a unitary category of dimension 2n\^{}2 with invertibles
C\_(n\^{}2), n simples of dimension sqrt(n), and the same cyclic
generalized-TY formulas. The two equivariantizations have dimension
4n\^{}2 and rank(n\^{}2+n+6)/2. The two rank-nine models above are the
n=3 members. This uniform construction does not classify all generalized
Tambara--Yamagami categories or their actions.

\hypertarget{an-order-four-cover-and-a-coherent-descended-action}{%
\section{An order-four cover and a coherent descended
action}\label{an-order-four-cover-and-a-coherent-descended-action}}

\hypertarget{result-and-outline}{%
\subsection{Result and outline}\label{result-and-outline}}

The ring R9\_2\_1254 is realized by a coherent C2 equivariantization of
a unitary generalized Moore--Read category of dimension 16. The new
action comes from a mixed order-four tensor functor on the dimension 32
covering category. Its square becomes tensor-isomorphic to the identity
only after central boson condensation. Requiring the cover functor to
have square identity would incorrectly discard this construction.

The certificate is
\texttt{mr\_\allowbreak{}order4\allowbreak{}\_\allowbreak{}descen\allowbreak{}t\_\allowbreak{}2\_\allowbreak{}3\_\allowbreak{}1\_\allowbreak{}4.\allowbreak{}json}:
it gives every root-of-unity tensorator scalar and every component of
the square isomorphism.
\texttt{mr\_\allowbreak{}order4\allowbreak{}\_\allowbreak{}descen\allowbreak{}t.\allowbreak{}py}
checks all required equations. The categorical argument below explains
why those equations suffice.

\hypertarget{the-existing-unitary-cover-and-its-central-algebra}{%
\subsection{The existing unitary cover and its central
algebra}\label{the-existing-unitary-cover-and-its-central-algebra}}

Use a unitary Ising category I with simples 1,psi,sigma, and
P=Vec(C16,Omega),

\begin{verbatim}
Omega(a,b,c)=zeta\^(2a floor((b+c)/16)), zeta=exp(2 pi i/16).
\end{verbatim}

Let B be the parity subcategory of I boxtimes P: its simples are(i,a),
with i=0,1,2 for1,psi,sigma, and a even for i=0,1 and odd for i=2. There
are24 simples and FPdim(B)=32. The invertible

\begin{verbatim}
g=(psi,8)
\end{verbatim}

has right half-braiding g tensor(i,a)-\textgreater(i,a)tensor g of
scalar

\begin{verbatim}
beta(i,a)=zeta\^(7a+e\_i),  (e\_0,e\_1,e\_2)=(0,8,12).
\end{verbatim}

It has order2 in Z(B), self-braiding1, and trivial associator on its
powers. Indeed the pointed square correction is zeta\^{}(2a), giving
beta\_pointed(a)\^{}2*zeta\^{}(2a)=1; the Ising central fermion squares
to1. The regular algebra A=1+g is commutative, separable and central.
Its multiplication g tensor g-\textgreater1 is the identity in the
chosen skeletal basis. Thus D=B\_A is a unitary fusion category of
dimension 16.

These are precisely the explicitly verified cover and central algebra in
\texttt{phase2\allowbreak{}/\allowbreak{}weak\_\allowbreak{}integr\allowbreak{}al/\allowbreak{}mr\_\allowbreak{}orbifo\allowbreak{}ld.\allowbreak{}md},
with the right-half-braiding orientation made explicit.

\hypertarget{the-order-four-mixed-tensor-functor}{%
\subsection{The order-four mixed tensor
functor}\label{the-order-four-mixed-tensor-functor}}

On simple objects define

\begin{verbatim}
F(i,a)=(i xor (a/2 mod 2),3a mod 16),   i=0,1;
F(2,a)=(2,3a mod 16).
\end{verbatim}

This preserves every fusion coefficient and the unit. The file lists the
scalar J\_F(a,b;c) in mu\_16 on every multiplicity-one fusion channel,
with tensorator orientation F(a)tensor F(b)-\textgreater F(a tensor b).
F is linear on each chosen elementary morphism and permutes channel
labels in the obvious way. Its tensor-functor equations are verified
entry by entry against the cover's actual F-symbols, not only the fusion
rules. All nonzero Ising F-symbol magnitudes are preserved: the only
nontrivial magnitude1/sqrt2 occurs for three sigma inputs, and F fixes
the Ising type of those inputs. Consequently checking the phase of every
nonzero entry checks the entire tensor-functor diagram.

The actual F-symbol exponent is the standard Ising sign times8, plus
\texttt{2 a fl\allowbreak{}oor((b\allowbreak{}+c)/\allowbreak{}16)} from
the charge factor. The script checks all 18432 nonzero associator
entries and the unit constraints. Since J is unitary and the object map
is a permutation, F is a unitary tensor autoequivalence.

F(g)=g. The certificate additionally verifies every central transport
identity

\begin{verbatim}
J\_F(g,a;ga)/J\_F(a,g;ga)=beta(Fa)/beta(a).
\end{verbatim}

Choose an algebra isomorphism phi:F(A)-\textgreater A which is identity
on1 and has scalar c=zeta32\^{}2 on g (equivalently c=zeta16). Here the
saved value is J\_F(g,g;1)=zeta16\^{}2, so c\^{}2=J\_F(g,g;1). This is
exactly the algebra multiplication compatibility. The central transport
equations make phi an isomorphism of central algebras. Therefore F
descends to a unitary tensor autoequivalence f of D.

\hypertarget{squaring-before-and-after-condensation}{%
\subsection{Squaring before and after
condensation}\label{squaring-before-and-after-condensation}}

On B, F\^{}2 is not the identity object map. It is

\begin{verbatim}
K(a)=g\^p(a) tensor a,
\end{verbatim}

where p(a)=0 for an even Ising object and1 for a sigma object. Equip K
with its canonical tensorator: move the central g factors to the left
using the right half-braiding, and multiply g\^{}2=1. More explicitly,
reassociate

\begin{verbatim}
(g\^p a)(g\^q b) -> g\^p(a(g\^q b))
  -> g\^p((a g\^q)b) -> g\^p((g\^q a)b)
  -> g\^p(g\^q(ab)) -> (g\^p g\^q)(ab) -> g\^(p+q)(ab).
\end{verbatim}

The middle interchange contributes beta(a)\^{}(-q). All other
contributions are specified cover associators. The script independently
verifies the resulting J\_K against all 18432 tensor-functor equations.

The certificate supplies a natural isomorphism d:F\^{}2-\textgreater K
on the common underlying simple objects. Its phases satisfy for every
channel a b-\textgreater c

\begin{verbatim}
d(a)d(b)/d(c)
  =J\_F(a,b;c) J\_F(Fa,Fb;Fc)/J\_K(a,b;c).
\end{verbatim}

This is exactly monoidality of d.~It also satisfies

\begin{verbatim}
d(g)=J\_F(g,g;1)=c\^2.
\end{verbatim}

Thus d intertwines the algebra map phi F(phi):F\^{}2(A)-\textgreater A
with the identity algebra map K(A)-\textgreater A. It descends to the
module categories.

On D, K is canonically tensor-isomorphic to the identity. On a free
module q(a), use q(g)\^{}p tensor q(a)-\textgreater q(a), where
q(g)-\textgreater1\_D is supplied by the regular algebra multiplication.
The tensorator of K was defined by precisely the central movements
appearing in this map. Its monoidality uses beta(g)=1 and the
group-algebra relation g\^{}2=1. Hence f\^{}2 is tensor-isomorphic to
Id\_D.

\hypertarget{coherent-c2-action-with-no-unresolved-obstruction}{%
\subsection{Coherent C2 action, with no unresolved
obstruction}\label{coherent-c2-action-with-no-unresolved-obstruction}}

For completeness there are two ways to finish coherence.

First, the canonical universal grading of D is C8: writing its simples
as u\^{}a(a in C8) and X\_r(r in C4), their degrees are2a and2r+1 modulo
8. Its adjoint subcategory has basis1,u\^{}4, as follows from X\_r
X\_r\^{}*=1+u\^{}4. Therefore Aut\_tensor(Id\_D)=Hom(C8,C*)=C8. The
functor f acts on this group by multiplication3. The obstruction to
upgrading an order-two tensor-equivalence class to a C2 action lies in

\begin{verbatim}
H\^3(C2,C8\_3)=ker(1+3)/(3-1)C8=2C8/2C8=0.
\end{verbatim}

Thus the action coherence can always be corrected. This is the standard
monoidal-action obstruction (a direct cyclic coherence calculation), not
an assertion that H4 alone supplies an extension.

Second, for the displayed certificate the correction is explicit. If
phases are written as exponents mod 16, replace d(a) by

\begin{verbatim}
d\_new(i,a)=d(i,a)-4a modulo 16.
\end{verbatim}

This multiplies d by the universal-grading character with parameter-2.
Before correction the C2 coherence defect is the parity character; after
correction it is identically1. In the saved bases the defect, written as
an exponent mod 32, is

\begin{verbatim}
2d(Fa)-2d(a)-p(a)(J\_F(g,g;1)-2J\_F(g,a;ga)).
\end{verbatim}

The script
\texttt{mr\_\allowbreak{}1254\_\allowbreak{}model.\allowbreak{}py}
checks that it is16 on odd objects and0 on even objects before
correction, and0 everywhere afterwards; the final phases still satisfy
all monoidality and algebra compatibility equations. Every correction is
unitary. This produces an actual coherent unitary C2 action on D.

\hypertarget{fusion-tensor-and-census-identification}{%
\subsection{Fusion tensor and census
identification}\label{fusion-tensor-and-census-identification}}

In D use representatives(1,2a) for u\^{}a, and(sigma,1+2r) for X\_r.
Then

\begin{verbatim}
u\^a u\^b=u\^(a+b),     u\^a X\_r=X\_r u\^a=X\_(r+a),
X\_r X\_s=u\^(r+s+1)+u\^(r+s+5).
\end{verbatim}

The descended involution acts by

\begin{verbatim}
u\^a ->u\^(-a),    X\_r ->X\_(3r+1).
\end{verbatim}

It fixes1,u\^{}4; it has three two-element pointed orbits, and two
two-element odd orbits. The equivariantization E=D\^{}C2 therefore has
rank 9 and dimensions 1\^{}4,2\^{}3,(2sqrt2)\^{}2. The odd orbits are
exchanged by duality, giving the feature missing from the earlier
model389.

For the fixed object u\^{}4, the certificate has
\texttt{d\_\allowbreak{}new(u\textasciicircum{}\allowbreak{}4)=14} and
\texttt{J\_\allowbreak{}F(u\textasciicircum{}4,\allowbreak{}u\textasciicircum{}4;1)\allowbreak{}=14}
modulo 16. Thus its equivariant lifts have square1, and the invertible
group is V4. A later action-coherence correction by a universal-grading
character cannot change this, because deg(u\^{}4)=0.

The full tensor is computed by the invariant-subspace formula. For a
fixed output the multiplicity is(dim+trace)/2; only fixed input objects
contribute to the trace. These inputs are1 andu\^{}4, whose
representatives are literally fixed in the cover, so the relevant phases
are exactly the saved J\_F and d\_new entries. For a free output orbit
use the dimension at one representative. This determines every fusion
coefficient.

\texttt{mr\_\allowbreak{}1254\_\allowbreak{}model.\allowbreak{}py}
verifies all ring axioms and matches all 729 coefficients of R9\_2\_1254
with source-to-model basis permutation

\begin{verbatim}
[0,8,7,5,6,1,2,4,3].
\end{verbatim}

\texttt{mr\_\allowbreak{}1254\_\allowbreak{}model.\allowbreak{}json}
includes the full rank 9 tensor, rank 12 quotient tensor, action,
corrected coherence phases and equivariant orbit labels. The model is
unitary by construction from unitary categories, a central unitary
separable algebra, and a coherent unitary finite-group action.

\hypertarget{scope-and-references}{%
\subsection{Scope and references}\label{scope-and-references}}

The proof establishes categorical existence for this fusion ring. It
does not classify all its categorifications or assert tensor-equivalence
novelty of the resulting category.

\begin{itemize}
\tightlist
\item
  S.Gelaki, D.Naidu, D.Nikshych, Centers of graded fusion categories,
  arXiv:0905.3117, Sections 2.4--2.5, for actions and central
  condensation.
\item
  J.Dong, S.Natale, H.Sun, A class of prime fusion categories of
  dimension 2\^{}N, New York J.Math.27(2021),141--163,
  \url{https://nyjm.albany.edu/j/2021/27-5p.pdf}, for the generalized
  Moore--Read setting. Existence and the required action here are
  specified directly.
\item
  P.Etingof, D.Nikshych, V.Ostrik, Fusion categories and homotopy
  theory, arXiv:0909.3140, for the obstruction-theoretic formulation of
  tensor actions; the relevant cyclic H3 calculation is supplied above.
\end{itemize}

\hypertarget{complete-group-theoretical-obstructions-in-dimensions-120-and-168}{%
\section{Complete group-theoretical obstructions in dimensions 120 and
168}\label{complete-group-theoretical-obstructions-in-dimensions-120-and-168}}

Scope. This argument excludes group-theoretical categorifications of
\texttt{R9\_\allowbreak{}2\_\allowbreak{}12},
\texttt{R9\_\allowbreak{}2\_\allowbreak{}959}, and
\texttt{R9\_\allowbreak{}2\_\allowbreak{}1965}. Together with ENO
Proposition 9.11 for the simple exotic ring
\texttt{R9\_\allowbreak{}2\_\allowbreak{}941}, all four integral
residual rings are outside the group-theoretical class. This does not
exclude arbitrary fusion categorifications. In particular, no theorem
identifying all integral categories of dimensions 120 or 168 as
group-theoretical is used.

The finite computations below are exact. A separate independent audit
passed:
\texttt{phase4\allowbreak{}/\allowbreak{}audit\_\allowbreak{}integr\allowbreak{}al\_\allowbreak{}gt/\allowbreak{}AUDIT.\allowbreak{}md}
and \texttt{audit\_\allowbreak{}verdic\allowbreak{}t.\allowbreak{}json}.
The accepted conclusion has group-theoretical scope only.

\hypertarget{an-exhaustive-projective-character-profile-reduction}{%
\subsection{An exhaustive projective-character profile
reduction}\label{an-exhaustive-projective-character-profile-reduction}}

Every group-theoretical category of dimension D has a presentation
\texttt{C(G,om\allowbreak{}ega,H,\allowbreak{}psi)} with \texttt{|G|=D},
\texttt{d psi=\allowbreak{}omega|\allowbreak{}H}. Its simples over a
double coset \texttt{HgH} are indexed by projective irreducible
representations of
\texttt{K\_\allowbreak{}g=H in\allowbreak{}tersec\allowbreak{}t gHg\textasciicircum{}\allowbreak{}\{-\allowbreak{}1\}},
of dimensions

\texttt{[H:K\_\allowbreak{}g] dim\allowbreak{}(pi)}.

The identity double coset supplies the ordinary category
\texttt{Rep(H)\allowbreak{}} as a fusion subcategory. A projective
irreducible degree is at most the square root of the order of its group,
so all simple dimensions are at most \textbar H\textbar. In particular,
\textbar H\textbar\textless6 cannot supply a dimension-six simple.

\texttt{all\_\allowbreak{}twist\_\allowbreak{}skelet\allowbreak{}on.\allowbreak{}g}
enumerates every group of orders 120 and 168 and every conjugacy class
of subgroups. It first checks that the ordinary character degrees of H
are a submultiset of the target type. For every remaining double-coset
intersection K, it constructs a Schur cover and partitions the
irreducible characters of the cover by their central character on the
kernel. These are precisely all projective degree profiles, for every
class in
\texttt{H\textasciicircum{}2(K,\allowbreak{}C\textasciicircum{}*)}. It
independently chooses any of these profiles at every double coset,
retaining those whose union is the target type. Independence is a
relaxation of global compatibility; it cannot omit any global
three-cocycle or relative two-cochain.

The complete output is:

\begin{longtable}[]{@{}rrrl@{}}
\toprule\noalign{}
D & Groups & Subgroup classes & Surviving pairs \\
\midrule\noalign{}
\endhead
\bottomrule\noalign{}
\endlastfoot
120 & 47 & 1431 & \texttt{(S5,S3\allowbreak{}xC2)} and
\texttt{(A5xC2\allowbreak{},S3xC2\allowbreak{})} \\
168 & 57 & 1751 & the two conjugacy classes
\texttt{(PSL(2\allowbreak{},7),S4\allowbreak{})} \\
\end{longtable}

Both dimension-120 pairs have intersections of orders
\texttt{12,2,4\allowbreak{}}, with the order-four intersection V4. Their
required projective degree packets are
\texttt{(1,1,1\allowbreak{},1,2,2\allowbreak{})}, \texttt{(6,6)}, and
\texttt{(6)} respectively. Thus the V4 intersection must carry its
nontrivial Schur class. Both dimension-168 pairs have intersections S4
and V4; the required packets are \texttt{(1,1,2\allowbreak{},3,3)} and
\texttt{(6,6,6\allowbreak{},6)}.

The scan uses GAP 4.12.1 with SmallGrp and the exact
character/Schur-cover algorithms. It is finite group classification plus
finite exact arithmetic, not a restricted three-cocycle ansatz. The
saved \texttt{.\allowbreak{}json} normalizes the GAP output; the
original \texttt{.\allowbreak{}jsonl} contains line-wrapped JSON
objects.

\hypertarget{immediate-consequence-for-r9_2_959}{%
\subsection{Immediate consequence for
R9\_2\_959}\label{immediate-consequence-for-r9_2_959}}

For both surviving dimension-120 pairs, H is \texttt{S3xC2}, whose
representation category has invertible group V4. It embeds as the
identity double-coset subcategory. But the invertible group of
R9\_2\_959 is C4. A based inclusion cannot embed V4 into C4. Thus no
group-theoretical categorification of R9\_2\_959 exists, independently
of omega and psi.

\hypertarget{all-relative-three-cocycles-for-the-s5-pair}{%
\subsection{All relative three-cocycles for the S5
pair}\label{all-relative-three-cocycles-for-the-s5-pair}}

Let \texttt{G=S5} and \texttt{H=S3xC\allowbreak{}2}. Every admissible
class
\texttt{a in H\allowbreak{}\textasciicircum{}3(G,C\allowbreak{}\textasciicircum{}*)}
restricts to zero on H. Restriction followed by corestriction is
multiplication by \texttt{[G:H]=\allowbreak{}10}, hence \texttt{10a=0}.

The five-primary part of
\texttt{H\textasciicircum{}3(S5\allowbreak{},C\textasciicircum{}*) \allowbreak{}= H\textasciicircum{}4(\allowbreak{}S5,Z)}
is zero: restriction to a Sylow C5 is injective on this part, and the
image must be fixed by the normalizer \texttt{C5:C4}. An order-four
automorphism of C5 acts on H\^{}4 by the square of a primitive
automorphism, namely -1 modulo five, whose fixed space is zero.
Therefore every admissible class has order at most two.

Every two-torsion class in H\^{}3(G,C\^{}*) lifts to H\^{}3(G,F2), by
the long exact sequence of
\texttt{1 -\allowbreak{}> mu2 \allowbreak{}-\allowbreak{}> C\textasciicircum{}* \allowbreak{}-\allowbreak{}-\allowbreak{}square\allowbreak{}-\allowbreak{}-\allowbreak{}> C\textasciicircum{}* \allowbreak{}-\allowbreak{}> 1}.
Let P be a Sylow D8 of S5. Its index is fifteen, hence odd. If a is a
two-torsion global class and b is any F2 lift of its restriction to P,
then the image of
\texttt{cor\_\allowbreak{}P\textasciicircum{}G(b)\allowbreak{}} in
H\^{}3(G,C\^{}*) is
\texttt{cor(re\allowbreak{}s(a))=\allowbreak{}15a=a}. Consequently
corestricting every F2 third cohomology class of P exhausts all
admissible global associator classes.

The normalized bar complex computes
\texttt{dim H\textasciicircum{}\allowbreak{}3(D8,F\allowbreak{}2)=4}:
343 normalized three-cochain variables, cocycle dimension 45, boundary
dimension 41. The program enumerates all sixteen classes. For right
coset representatives \texttt{t} of \texttt{P\textbackslash{}G}, define
\texttt{h(t,g)\allowbreak{}} and \texttt{t.\allowbreak{}g} by

\texttt{t g = \allowbreak{}h(t,g)\allowbreak{} (t.\allowbreak{}g)}.

The explicit corestricted sign cocycle is

\texttt{w(g1,g\allowbreak{}2,g3) \allowbreak{}= sum\_\allowbreak{}t b(h(\allowbreak{}t,g1),\allowbreak{}h(t.\allowbreak{}g1,g2)\allowbreak{},h(t.\allowbreak{}g1g2,g\allowbreak{}3)) mo\allowbreak{}d 2}.

This is the usual inhomogeneous cochain transfer, so it takes exact
cocycles to exact cocycles. The program solves
\texttt{d psi=\allowbreak{}(-\allowbreak{}1)\textasciicircum{}w|H\allowbreak{}}
exactly in eighth roots of unity. Completeness of this finite primitive
test is explained below.

Eight of the sixteen representatives admit a relative trivialization.
Four yield rank twelve. The other four yield rank nine and match **all
729 coefficients of R9\_2\_92**, with explicit based isomorphisms; none
matches R9\_2\_12. Reusing the independently audited twisted bimodule
character formula is legitimate here: only its normalized cocycle and
relative cochain inputs change. Every returned tensor is separately
checked for the unit and every associativity equation: 6,561 at rank
nine and 20,736 at rank twelve.

There is no missing choice of psi.
\texttt{H\textasciicircum{}2(H,\allowbreak{}C\textasciicircum{}*)=C\allowbreak{}2},
and the nontrivial class extends from the Schur cocycle of S5: lifted
disjoint transpositions have nontrivial commutator. Multiplying psi by
that class is induced by an identity-on-objects tensor autoequivalence
of \texttt{Vec\_\allowbreak{}G\textasciicircum{}omeg\allowbreak{}a}, so
it does not change the based ring. This extension was independently
checked in phase three's
\texttt{S5\_\allowbreak{}RELATI\allowbreak{}VE\_\allowbreak{}TWIST.\allowbreak{}md}
and audit.

Thus the entire S5 pair, with all relative cocycles, gives no model for
R9\_2\_12.

\hypertarget{exact-primitive-solver-and-its-bound}{%
\subsubsection{Exact primitive solver and its
bound}\label{exact-primitive-solver-and-its-bound}}

Since \texttt{H\textasciicircum{}2(H,\allowbreak{}C\textasciicircum{}*)}
has exponent two, a sign-valued three-cocycle on H that is a
C\^{}*-coboundary has a fourth-root-valued primitive. Indeed, if
\texttt{d psi=\allowbreak{}w}, then \texttt{psi\textasciicircum{}2} is a
two-cocycle. Hence \texttt{psi\textasciicircum{}4=\allowbreak{}d f} for
a one-cochain f.~Choose a fourth root b of f, and replace psi by
\texttt{psi/\allowbreak{}(d b)}; its fourth power is one. Thus eighth
roots are more than sufficient.

The congruence solver lifts successively modulo two, four and eight. At
each lower stage it enumerates all possible H\^{}2(H,F2) classes, whose
basis is
\texttt{x\textasciicircum{}2,xy\allowbreak{},y\textasciicircum{}2}, for
the two abelianization characters of \texttt{H=S3xC\allowbreak{}2}. This
follows also from the normal odd-order C3 subgroup:
\texttt{H\textasciicircum{}*(H,\allowbreak{}F2)=H\textasciicircum{}\allowbreak{}*(C2xC\allowbreak{}2,F2)}.
Two solutions differing by an F2 two-coboundary have the same
liftability, since that change can be lifted to an integral
one-coboundary adjustment. Exact Gaussian elimination handles the
inhomogeneous linear equations and tests every consistency row. The
final primitive is checked against every normalized triple of H. Thus a
failed test is an exact no-primitive conclusion, not an optimizer
timeout.

\hypertarget{the-a5xc2-pair-including-the-odd-primary-part}{%
\subsection{The A5xC2 pair, including the odd-primary
part}\label{the-a5xc2-pair-including-the-odd-primary-part}}

Again the global restriction kernel is killed by ten, so an admissible
class is a sum of a two-torsion and a five-torsion class. We handle the
two-primary part by the same exhaustive transfer construction, using a
Sylow group P=C2\^{}3 of index fifteen. The normalized bar complex has
cocycle dimension 49 and boundary dimension 39, giving ten F2 third
cohomology generators.

For efficiency, let n be an order-three element of the normalizer of P.
Corestriction is unchanged in cohomology by conjugating a seed by n.
Thus corestricting b is cohomologous to corestricting
\texttt{b+n*b+\allowbreak{}n\textasciicircum{}2*b}, since three is one
modulo two. The span of these norm classes has dimension four. Exact
reduction modulo the bar boundaries therefore gives sixteen
representatives covering the image of **all 1,024** original F2
cohomology classes. No global class is discarded.

Only two of these sixteen representatives admit a primitive on H. For
each, both subgroup Schur choices \texttt{psi} and
\texttt{psi*(-\allowbreak{}1)\textasciicircum{}\{x \allowbreak{}cup y\}\allowbreak{}}
are explicitly tested. In all four cases the V4 intersection multiplier
is trivial and the resulting category has rank twelve. Thus no
two-primary relative class gives the required rank-nine profile.

The five-primary part cannot change this conclusion. Relative pairs
\texttt{(omega\allowbreak{},psi)} form relative cohomology, and the
multiplier on a double-coset intersection K is a homomorphism from that
relative cohomology group to
\texttt{H\textasciicircum{}2(K,\allowbreak{}C\textasciicircum{}*)}. The
five-primary part of the relative group maps isomorphically to the
five-primary global restriction kernel, because
\texttt{H\textasciicircum{}2(H,\allowbreak{}C\textasciicircum{}*)} has
no five-torsion. Its image in the order-two Schur multiplier of V4 is
therefore zero. Equivalently, multiplying the relative data by any
order-five relative class does not change the V4 projective degree
profile. The necessary rank-nine profile requires that multiplier to be
nontrivial. Hence no admissible cocycles for
\texttt{(A5xC2\allowbreak{},S3xC2\allowbreak{})} realize R9\_2\_12.

Combining Sections 1, 3 and 4 proves the group-theoretical exclusion for
R9\_2\_12.

\hypertarget{the-dimension-168-pair}{%
\subsection{The dimension-168 pair}\label{the-dimension-168-pair}}

The all-twist profile reduction leaves only
\texttt{G=PSL(\allowbreak{}2,7)}, H=S4, for the two subgroup conjugacy
classes. The following completes the old bounded argument by showing
that these are now the *only* presentations that need to be considered.

Restriction
\texttt{H\textasciicircum{}3(G,\allowbreak{}C\textasciicircum{}*) -\allowbreak{}> H\textasciicircum{}3(\allowbreak{}H,C\textasciicircum{}*)\allowbreak{}}
is injective. Its kernel is killed by index seven. But G has no
seven-primary H\^{}4: restriction to its Sylow C7 is injective on the
seven-primary part, and the normalizer C7:C3 acts on H\^{}4(C7,Z)=C7 by
multiplication by 2 and 4, with zero fixed part. Thus any admissible
omega is globally trivializable.

After trivializing omega,
\texttt{H\textasciicircum{}2(S4\allowbreak{},C\textasciicircum{}*)=\allowbreak{}C2},
and both classes extend from
\texttt{H\textasciicircum{}2(PS\allowbreak{}L(2,7)\allowbreak{},C\textasciicircum{}*)=\allowbreak{}C2}.
Nontriviality of restriction is witnessed by a commuting pair in S4
whose lifts to the Schur double cover fail to commute. The prior exact
checker tests this for both subgroup classes:
\texttt{phase3\allowbreak{}/\allowbreak{}group\_\allowbreak{}models\allowbreak{}/\allowbreak{}psl168\allowbreak{}\_\allowbreak{}relati\allowbreak{}ve\_\allowbreak{}bound.\allowbreak{}g}
and its saved log. Therefore psi can also be removed by a global
two-cocycle autoequivalence. Both presentations reduce to the untwisted
category, whose full tensor is R9\_2\_1557, not R9\_2\_1965. The based
nonisomorphism is checked again.

\hypertarget{fermion-parity-and-solvable-descent-in-dimension-120}{%
\section{Fermion parity and solvable descent in dimension
120}\label{fermion-parity-and-solvable-descent-in-dimension-120}}

\hypertarget{a-fermion-parity-criterion}{%
\subsection{A fermion-parity
criterion}\label{a-fermion-parity-criterion}}

Proposition. Let C be an integral fusion category with commutative
Grothendieck ring, trivial universal grading, and dimension N divisible
by four. Write the induction of the unit into Z(C) as a sum of distinct
simple objects. Suppose exactly one nonunit summand has odd dimension,
and that its dimension is a positive power of an odd prime. Then Z(C)
contains a nontrivial Tannakian subcategory.

Proof. Integrality gives a canonical positive spherical structure and
makes Z(C) an integral nondegenerate braided category. The
odd-prime-power summand of I(1) and ENO Corollary 7.2 produce a
nontrivial symmetric fusion subcategory. If that symmetric category has
dimension larger than two, Deligne's theorem gives a nontrivial
Tannakian subcategory. A symmetric category of dimension two is either
Tannakian or sVec. It remains to handle sVec, generated by a fermion chi
in Z(C).

Here is a short proof of the required dimension count. By Ostrik
Proposition 2.8, for the unnormalized modular S-matrix,

\begin{verbatim}
S [I(1)] = N [I(1)].
\end{verbatim}

The fermion chi has twist -1, whereas every summand of A=I(1) has twist
+1 (Ostrik Theorem 2.5); hence chi is absent from A. The chi entry of
this identity is zero. For each simple Z in A, monodromy with the
order-two invertible chi is a scalar epsilon\_Z in \{+1,-1\}, and
\texttt{S\_\allowbreak{}\{chi,Z\allowbreak{}\}=epsi\allowbreak{}lon\_\allowbreak{}Z FPdi\allowbreak{}m(Z)}.
Thus the total dimension of the constituents of A centralizing chi
equals that of the other constituents, namely N/2. This uses no braiding
on C.

For completeness the same dimension count has the following independent
relative-center derivation. Let F:Z(C)-\textgreater C be the forgetful
functor and g=F(chi). Because the universal grading of C is trivial, the
forgetful homomorphism on invertible objects is injective: its kernel is
the group of tensor automorphisms of Id\_C, namely the characters of the
universal grading group. Thus g is a nontrivial order-two invertible,
and D= has dimension two.

Use the relative-center construction from the proof of ENO Lemma 9.17.
The forgetful functor Z(C)-\textgreater Z\_D(C), which remembers the
half-braiding only against D, is dominant. Its right adjoint on the unit
gives a connected algebra B contained in A, of dimension N/FPdim(D)=N/2.
This follows either from that lemma's dimension calculation or from the
dominant-functor index formula and FPdim(Z\_D(C))=FPdim(C) FPdim(D).

We claim that B is precisely the sum of those constituents of A that
centralize chi. Here is the explicit morphism check, so this
identification does not require a fixed-algebra theorem or a choice of
associator. For a central simple Z occurring in A, take a nonzero map
f:F(Z)-\textgreater1. It is a morphism to the unit in the relative
center precisely when

\begin{verbatim}
(id\_g tensor f) c\_{Z,g} = f tensor id\_g.
\end{verbatim}

The naturality of chi's half-braiding gives

\begin{verbatim}
(f tensor id\_g) c\_{g,Z} = id\_g tensor f.
\end{verbatim}

Consequently the relative-center condition is equivalent to

\begin{verbatim}
(f tensor id\_g) (c\_{g,Z} c\_{Z,g}) = f tensor id\_g.
\end{verbatim}

In Z(C), Z tensor chi is simple, so this monodromy is scalar. Since chi
has order two, its monodromy eigenvalue is +1 or -1. The displayed
condition therefore holds exactly when Z centralizes chi. This proves
the claim, including multiplicities (which are one because the based
ring is commutative).

The dimension of B is N/2, which is even. B contains the unit, and every
nonunit summand of A except the specified odd-prime-power object has
even dimension. Parity therefore forces that odd object to occur in B.
It lies in the centralizer E=' inside Z(C). By nondegeneracy of Z(C),
the Muger center of E is exactly =sVec, so E is slightly degenerate. ENO
Proposition 7.4 says that a slightly degenerate integral braided
category containing a simple object of odd-prime-power dimension has a
nontrivial Tannakian subcategory. This proves the proposition. QED.

The argument explicitly treats a possible central fermion. It does not
infer a Tannakian category merely from an unspecified symmetric
category.

\hypertarget{application-to-both-dimension-120-tensors}{%
\subsection{Application to both dimension-120
tensors}\label{application-to-both-dimension-120-tensors}}

Both input tensors have dimensions
\texttt{(1\textasciicircum{}4,2\allowbreak{}\textasciicircum{}2,6\textasciicircum{}3\allowbreak{})},
are commutative, and have trivial universal grading. Their formal
codegrees, with multiplicity, are

\begin{verbatim}
120, 24, 20, 30, 30, 6, 6, 4, 4.
\end{verbatim}

Ostrik's induction theorem (Theorem 2.13 of arXiv:1309.4822) gives a
multiplicity-free decomposition of A=I(1) with simple dimensions

\begin{verbatim}
1, 5, 6, 4, 4, 20, 20, 30, 30.
\end{verbatim}

There is exactly one odd nonunit summand, of dimension five. The
proposition therefore forces a nontrivial Tannakian subcategory of Z(C).

Li Dai's Lemma 3.1 then implies that C is weakly group-theoretical. Its
full hypothesis is exactly: an integral category of dimension 120 whose
center contains a nontrivial Tannakian subcategory. For clarity, its
reduction is as follows. If the Tannaka group's order is a power of two,
it contains a quotient C2, and the corresponding Rep(C2) in the center
gives an extension or equivariantization of a dimension-60 category. All
dimension-60 categories are weakly group-theoretical by ENO Theorem
9.16. If the Tannaka group's order is divisible by three or five, the
neutral component of the de-equivariantized center has dimension
\texttt{(120/\allowbreak{}|G|)\textasciicircum{}2\allowbreak{}}, with at
most two prime divisors, and is solvable by ENO Theorem 1.6. The
extension/equivariantization and center characterizations of weak
group-theoreticality finish the reduction.

New conclusion: every hypothetical categorification of either
\texttt{R9\_\allowbreak{}2\_\allowbreak{}12} or
\texttt{R9\_\allowbreak{}2\_\allowbreak{}959} is weakly
group-theoretical. This conclusion applies to arbitrary complex fusion
categories and is not a model construction.

\hypertarget{an-auxiliary-dimension-30-obstruction}{%
\subsection{An auxiliary dimension-30
obstruction}\label{an-auxiliary-dimension-30-obstruction}}

Lemma. An integral fusion category of dimension 30 has no simple object
of dimension three.

Proof. ENO Theorem 9.2 makes every integral fusion category of
square-free dimension 30 group-theoretical. Every group G of order 30
has a normal cyclic subgroup P of order 15. One elementary proof is: if
its Sylow-five subgroup were not normal, there would be 24 nonidentity
order-five elements. Counting forces the Sylow-three subgroup to be
normal. The inverse image of the normal Sylow-five subgroup in the
order-ten quotient is then a normal subgroup of order 15, necessarily
cyclic, making its Sylow-five subgroup normal after all. Thus Sylow five
is normal. In the order-six quotient, the Sylow-three subgroup is
normal; its inverse image is the desired cyclic subgroup of order 15.

In a group-theoretical presentation C(G,omega,H,psi), simples over HgH
are indexed by projective irreducible representations of K=H intersect
gHg\^{}\{-1\}, with dimension
\texttt{[H:K] \allowbreak{}dim(pi\allowbreak{})}. Every such K has
square-free order, so H\^{}2(K,C*)=0: restriction to a Sylow subgroup is
injective on the corresponding primary part by
restriction/corestriction, while every Sylow subgroup is cyclic and has
trivial complex Schur multiplier. Thus all projective degrees coincide
with ordinary character degrees.

Set Q=H intersect P. It is cyclic and, being a subgroup of the cyclic
normal subgroup P, is normal in G. We have {[}H:Q{]}\textless=2, and Q
lies in every intersection K. Hence {[}H:K{]}\textless=2. If
{[}H:K{]}=2, K=Q is cyclic, so every projective irreducible of K has
dimension one. If {[}H:K{]}=1, K=H has a cyclic normal subgroup of index
at most two. Restriction of an ordinary irreducible representation to
that cyclic subgroup is a sum of characters. Clifford theory gives
irreducible dimensions at most two: an orbit of size two induces a
two-dimensional representation, while an invariant character extends
across the cyclic quotient and gives one-dimensional representations.
Thus all resulting simple dimensions are one or two. In particular no
dimension-three simple occurs. QED.

This argument allows arbitrary three-cocycles and subgroup cochains; an
untwisted character-table search alone would not suffice. The elementary
group-of-order-30 fact can equivalently be checked on the four groups of
order 30. It is also a consequence of the classification of groups of
square-free order at this particular order.

\hypertarget{no-solvable-category-has-dimension-type-1636}{%
\subsection{No solvable category has dimension type
(1\^{}6,3\^{}6)}\label{no-solvable-category-has-dimension-type-1636}}

Theorem. A fusion category D with six invertibles and six simple objects
of dimension three is not solvable.

Proof. Its dimension is 60. For each dimension-three simple X, G{[}X{]}
has order dividing both \textbar G(D)\textbar=6 and FPdim(X)\^{}2=9. In
the decomposition of X tensor X*, all noninvertible summands have
dimensions divisible by three, so \textbar G{[}X{]}\textbar{} is
divisible by three. Hence \textbar G{[}X{]}\textbar=3. The group of six
invertibles has a unique subgroup H of order three, so this same H fixes
all six noninvertible objects. It lies in D\_ad. The remainder of X
tensor X* has dimension six, so D\_ad is nonpointed.

Consider a nontrivial faithful grading with group order u. Each
component has dimension 60/u. The neutral component contains H and at
least one three-dimensional object. A component with a invertibles and b
three-dimensional simples has dimension a+9b. There are six invertibles
and six three-dimensional simples in all. These constraints leave only
u=2, with both components of type
\texttt{(1\textasciicircum{}3,3\allowbreak{}\textasciicircum{}3)}. This
elementary enumeration is replayed exhaustively by the verifier. The
neutral category would have dimension 30 and a three-dimensional simple,
contrary to Section 3. Thus U(D)=1.

If D were solvable, the final step of a cyclic-prime extension/
equivariantization construction could not be an extension. It must be an
equivariantization by C\_p.~Its central Rep(C\_p) embeds into D and
supplies p invertibles, so p=2 or p=3.

For p=2, its nontrivial character acts freely on the six dimension-three
objects: their stabilizers under G(D) have order three. It also acts
freely on all six invertibles. De-equivariantization has type
\texttt{(1\textasciicircum{}3,3\allowbreak{}\textasciicircum{}3)} and
dimension 30, again impossible.

For p=3, its character subgroup is H and fixes all six noninvertible
objects. The six invertibles form two free orbits. De-equivariantization
has dimension 20 and is pointed: each fixed three-dimensional simple
splits into three invertibles, and each free orbit of invertibles gives
one invertible. To justify the splitting, the endomorphism algebra of
the free module of a fixed simple is a twisted group algebra of C3;
H\^{}2(C3,C*)=0, so it has three one-dimensional simple modules. Each
has underlying object X and normalized dimension 3/3=1.

Write this pointed category as Vec\_K\^{}omega, \textbar K\textbar=20.
The equivariantizing C3 acts on K by group automorphisms. Its invariant
invertibles each contribute exactly three invertibles in the
equivariantization, again because H\^{}2(C3,C*)=0. Six invertibles of D
would require exactly two fixed elements of K. But K has a unique, hence
characteristic, Sylow-five subgroup, and the C3 action is trivial there
because Aut(C5)=C4. There are at least five fixed elements, a
contradiction. QED.

This theorem does not exclude categories of that dimension type in
general: the group-theoretical category C(A5,1,C3,1) supplies that type.
The conclusion is nonsolvability, which is the needed restriction.

\hypertarget{no-solvable-category-has-dimension-type-122326}{%
\subsection{No solvable category has dimension type
(1\^{}2,2,3\^{}2,6)}\label{no-solvable-category-has-dimension-type-122326}}

Theorem. A fusion category D of this type is not solvable.

Proof. A nontrivial faithful grading of a dimension-60 fusion category
has components of dimension at most 30. A simple object of dimension six
alone contributes 36 to its component, which is impossible. Hence
U(D)=1. If D were solvable, its last cyclic-prime step would be an
C2-equivariantization, since it has exactly two invertibles. Let g be
the nontrivial underlying character. It fixes the unique objects of
dimensions two and six. It exchanges the two dimension-three simples:
the stabilizer of a dimension-three simple has order dividing nine, so
cannot contain an involution. De-equivariantization therefore has three
invertibles (one from the pair of original invertibles and two from the
fixed dimension-two object) and three dimension-three simples (one from
their pair, two from the fixed dimension-six object). This is the
forbidden dimension-30 type
\texttt{(1\textasciicircum{}3,3\allowbreak{}\textasciicircum{}3)} of
Section 3. QED.

\hypertarget{central-involution-descents-of-the-two-source-rings}{%
\subsection{Central involution descents of the two source
rings}\label{central-involution-descents-of-the-two-source-rings}}

For a hypothetical central Tannakian Rep(C2), the underlying involution
must be one of the following source labels (zero-based):

\begin{longtable}[]{@{}lrl@{}}
\toprule\noalign{}
Ring & Involution & Dimension type after de-equivariantization \\
\midrule\noalign{}
\endhead
\bottomrule\noalign{}
\endlastfoot
R9\_2\_12 & 1 or 2 &
\texttt{(1\textasciicircum{}2,2\allowbreak{},3\textasciicircum{}2,6\allowbreak{})} \\
R9\_2\_12 & 7 &
\texttt{(1\textasciicircum{}6,3\allowbreak{}\textasciicircum{}6)} \\
R9\_2\_959 & 8 &
\texttt{(1\textasciicircum{}6,3\allowbreak{}\textasciicircum{}6)} \\
\end{longtable}

The verifier checks every orbit of multiplication by each involution. A
free pair of d-dimensional objects gives one d-dimensional object below;
a fixed d-dimensional object splits into two objects of dimension d/2.
This follows from the algebra 1+g and H\^{}2(C2,C*)=0, with all
categorical premises stated. Sections 4 and 5 rule out both dimension
types whenever the descending category is required to be solvable.

\hypertarget{nonsolvable-weakly-group-theoretical-categories-of-dimension-120}{%
\subsection{Nonsolvable weakly group-theoretical categories of dimension
120}\label{nonsolvable-weakly-group-theoretical-categories-of-dimension-120}}

Lemma. Every nonsolvable weakly group-theoretical fusion category of
dimension 120 is group-theoretical.

Proof. By definition it is Morita equivalent to a nilpotent category N,
built from Vec by successive faithful extensions by finite groups G\_i.
The product of the orders \textbar G\_i\textbar{} is 120. If every G\_i
is solvable, refine each extension by a composition series into
cyclic-prime extensions. Then N, and hence the original category, is
solvable, a contradiction.

Consequently one grading group is nonsolvable. A nonsolvable finite
group whose order divides 120 has a nonabelian composition factor A5:
the only possible proper nonsolvable order is 60, and its group is A5.
If a grading group has order 120, every other stage is trivial and N is
pointed. Otherwise an A5 extension uses a factor 60, leaving only a
single factor two. Up to deleting trivial stages, there are two possible
orders.

First suppose Vec is extended by A5 and then by C2. The intermediate
category is P=Vec\_\{A5\}\^{}omega. Every invertible P-bimodule
category, as a left P-module category, has pointed dual. The
classification of such module categories (Naidu--Nikshych, Example 2.2)
requires a normal abelian subgroup H of A5. Its only possibility is H=1,
so the left module is the regular P-module, with 60 simples permuted
transitively by A5. Each component of any extension of P is such a
bimodule, and has FP dimension 60. Its 60 equally dimensional simples
therefore all have dimension one. The extension N is pointed, for every
omega and all coherent extension data.

Second suppose Vec is extended by C2 and then by A5. The intermediate
category is P=Vec\_\{C2\}\^{}omega. Its center is pointed of rank four:
its simples are labeled by the two elements of C2 and projective
irreducible characters of C2, which are all one-dimensional since
H\^{}2(C2,C*)=0. The group BrPic(P) is the group of braided tensor
autoequivalences of Z(P), and embeds into the permutations of its three
nonunit objects. For clarity, faithfulness follows from the
classification of pointed braided categories by metric groups: a braided
tensor autoequivalence fixing all simple objects is tensor isomorphic to
the identity. Thus BrPic(P) is a subgroup of S3 and is solvable. The
homomorphism \texttt{A5 -\allowbreak{}> BrPi\allowbreak{}c(P)} of
extension components is trivial, since A5 is simple nonabelian. Every
component is the regular P-bimodule and has two simple objects. They
have equal FP dimension and their squared dimensions sum to FPdim(P)=2;
both are invertible. Again N is pointed.

In all cases N is pointed, so the original category is
group-theoretical. QED. This is a dimension-120 statement with the
nonsolvability hypothesis; it does not assert that all dimension-120
categories are group-theoretical.

\hypertarget{general-exclusion-of-r9_2_12-and-r9_2_959}{%
\subsection{General exclusion of R9\_2\_12 and
R9\_2\_959}\label{general-exclusion-of-r9_2_12-and-r9_2_959}}

Theorem. Neither ring admits a complex fusion categorification.

Proof. Suppose C categorifies either ring. Section 2 proves that C is
weakly group-theoretical. The independently audited phase-four theorem
\texttt{phase4\allowbreak{}/\allowbreak{}integr\allowbreak{}al/\allowbreak{}GROUP\_\allowbreak{}THEORE\allowbreak{}TICAL\_\allowbreak{}OBSTRU\allowbreak{}CTIONS\allowbreak{}.\allowbreak{}md}
excludes every group-theoretical categorification, including all
associators and subgroup cochains. Section 7 consequently forces C to be
solvable.

Its universal grading is trivial by the exact input tensor. The final
step of a cyclic-prime construction of C must therefore be an
equivariantization. Since C has four invertibles, the prime is two, and
the canonical central Rep(C2) has one of the underlying involutions
listed in Section 6. Its de-equivariantization D is solvable, being the
preceding category in this construction. It has one of the two dimension
types excluded for solvable categories by Sections 4 and 5. This is a
contradiction. QED.

Every reduction concerns an arbitrary hypothetical categorification. The
proof does not assume a fiber functor, braiding, unitarity, or an
uncited classification of all dimension-120 categories. It uses the
phase-four all-cocycle GT theorem only within its established scope.

\hypertarget{references-1}{%
\subsection{References}\label{references-1}}

* Etingof--Nikshych--Ostrik, *Weakly group-theoretical and solvable
fusion categories*, arXiv:0809.3031v2. Corollary 7.2; Proposition 7.4;
Propositions 2.9, 2.10 and 4.4; Theorems 9.2 and 9.16; proof of Lemma
9.17 (relative-center algebra). \url{https://arxiv.org/pdf/0809.3031} *
Ostrik, *Pivotal fusion categories of rank 3*, arXiv:1309.4822, Theorem
2.5, Proposition 2.8, and Theorem 2.13 (valid for arbitrary spherical
fusion categories). \url{https://arxiv.org/pdf/1309.4822} * Dai,
*Frobenius property for fusion categories of dimension 120*, Revista de
la Union Matematica Argentina 68(2) (2025), 395--403, Lemma 3.1. DOI
10.33044/revuma.4390.
\url{https://inmabb.criba.edu.ar/revuma/pdf/v68n2/v68n2a03.pdf}

* Naidu--Nikshych, *Lagrangian subcategories and braided tensor
equivalences of twisted quantum doubles of finite groups*,
arXiv:0705.0665, Example 2.2 (pointed-dual module categories).
\url{https://arxiv.org/pdf/0705.0665} *
Etingof--Gelaki--Nikshych--Ostrik, *Tensor Categories*, Sections 8.4
(pointed braided categories and metric groups), specifically Theorem
8.4.12, and Section 8.27.2 (Brauer--Picard groups and centers).
\url{https://math.mit.edu/~etingof/egnobookfinal.pdf}

The lemmas above are combinations of these primary results and explicit
arguments; no claim of literature priority is made. Final status and any
further phase-five inputs are recorded separately in REPORT.md.

\hypertarget{odd-central-dimensions-and-the-dimension-504-ring}{%
\section{Odd central dimensions and the dimension-504
ring}\label{odd-central-dimensions-and-the-dimension-504-ring}}

Theorem. The fusion ring \texttt{R9\_\allowbreak{}2\_\allowbreak{}941}
admits no fusion categorification over \texttt{C}. This conclusion uses
no unitarity assumption. The ring is exactly the rank-nine,
dimension-504 ring in Alekseyev--Bruns--Dong--Palcoux,
arXiv:2507.07023v2, §1.5, Question 1.14; the previously verified tensor
isomorphism is recalled below.

\hypertarget{notation-and-categorical-input}{%
\subsection{Notation and categorical
input}\label{notation-and-categorical-input}}

Use source labels \texttt{X0=1}, \texttt{Y=X2}, \texttt{B1=X1},
\texttt{B2=X3}, \texttt{B3=X4}, \texttt{B4=X5}, and \texttt{C1=X6},
\texttt{C2=X7}, \texttt{C3=X8}. Their dimensions are respectively
\texttt{1,8,7,\allowbreak{}7,7,7,\allowbreak{}9,9,9}, and total
dimension is \texttt{D=504}. Suppose a categorification C exists. Since
C is integral, it has the canonical positive spherical structure. Its
center Z(C) is an integral modular category of dimension
\texttt{D\textasciicircum{}2}. Write F for forgetful functor and I for
its right adjoint. For a central simple Z write

\texttt{F(Z) =\allowbreak{} t*1 +\allowbreak{} a*Y +\allowbreak{} sum\_\allowbreak{}j b\_\allowbreak{}j*B\_\allowbreak{}j + su\allowbreak{}m\_\allowbreak{}l c\_\allowbreak{}l*C\_\allowbreak{}l},

and let \texttt{z=dim(\allowbreak{}Z)},
\texttt{b=sum\_\allowbreak{}j b\_\allowbreak{}j},
\texttt{c=sum\_\allowbreak{}l c\_\allowbreak{}l}. Thus
\texttt{z=t+8a\allowbreak{}+7b+9c\allowbreak{}}.

The following standard facts will be used:

\begin{enumerate}
\def\labelenumi{\arabic{enumi}.}
\item
  \texttt{z | D}, because modular square-divisibility makes
  \texttt{D\textasciicircum{}2/\allowbreak{}z\textasciicircum{}2} an
  integer, and \texttt{D,z} are positive integers.
\item
  Since K0(C) is commutative, I(1) is multiplicity-free. Its central
  simples A\_chi are indexed by the nine scalar characters, with
  \texttt{dim(A\_\allowbreak{}chi)=D\allowbreak{}/\allowbreak{}f\_\allowbreak{}chi}.
  Thus \texttt{t} is zero or one. The central dimensions in I(1), with
  multiplicity, are
  \texttt{1,63,5\allowbreak{}6,56,5\allowbreak{}6,56,7\allowbreak{}2,72,7\allowbreak{}2}.
\item
  If chi is a scalar character, the central character identity gives
  \texttt{S\_\allowbreak{}(A\_\allowbreak{}chi,Z)\allowbreak{}=dim(A\allowbreak{}\_\allowbreak{}chi)*c\allowbreak{}hi(F(Z\allowbreak{}))}.
  By symmetry of S and the Verlinde character formula,
  \texttt{dim(A\_\allowbreak{}chi)*c\allowbreak{}hi(F(Z\allowbreak{}))/\allowbreak{}z}
  is an algebraic integer.
\item
  Set
  \texttt{G=sum\_\allowbreak{}i N\_\allowbreak{}i N\_\allowbreak{}i\textasciicircum{}T}.
  The induction Gram identity gives
  \texttt{G=sum\_\allowbreak{}Z v\_\allowbreak{}Z v\_\allowbreak{}Z\textasciicircum{}T},
  where v\_Z is the restriction column. Consequently
  \texttt{v\_\allowbreak{}Z\textasciicircum{}T G\textasciicircum{}\allowbreak{}(-\allowbreak{}1) v\_\allowbreak{}Z <=1}.
  Orthogonality of the real scalar characters gives

  \texttt{v\_\allowbreak{}Z\textasciicircum{}T G\textasciicircum{}\allowbreak{}(-\allowbreak{}1) v\_\allowbreak{}Z = su\allowbreak{}m\_\allowbreak{}chi ch\allowbreak{}i(F(Z)\allowbreak{})\textasciicircum{}2/\allowbreak{}f\_\allowbreak{}chi\textasciicircum{}2 \allowbreak{}<=1}.
  (1)

  This is an ordinary positive semidefinite matrix identity and requires
  no C*-structure on C.
\item
  Frobenius reciprocity and dimensions give

  \texttt{sum\_\allowbreak{}Z z*[F\allowbreak{}(Z):X]\allowbreak{} = dim\allowbreak{}(I(X))\allowbreak{} = D*d\allowbreak{}im(X)}.
  (2)
\end{enumerate}

For the identity in item 3, see Deshpande, arXiv:1611.04486, Proposition
2.8 and §2.7, especially (36) and the description immediately preceding
it. For item 2 see Ostrik, arXiv:1309.4822, Theorem 2.13.

\hypertarget{exact-character-data}{%
\subsection{Exact character data}\label{exact-character-data}}

The dimension character has formal codegree 504. One rational character
of codegree 8 is

\texttt{chi\_\allowbreak{}8(F(Z)\allowbreak{}) = t-\allowbreak{}b+c}.

Four rational characters of codegree 9 are

\texttt{chi\_\allowbreak{}9,j(F(\allowbreak{}Z)) = \allowbreak{}t+b-\allowbreak{}a-\allowbreak{}3*b\_\allowbreak{}j}
(\texttt{1<=j<=\allowbreak{}4}).

The other three characters have codegree 7. Put
\texttt{alpha=\allowbreak{}2*cos(\allowbreak{}2*pi/\allowbreak{}7)},
\texttt{beta=a\allowbreak{}lpha\textasciicircum{}2\allowbreak{}-\allowbreak{}2},
\texttt{gamma=\allowbreak{}1-\allowbreak{}alpha-\allowbreak{}alpha\textasciicircum{}\allowbreak{}2}.
These are the roots of
\texttt{x\textasciicircum{}3+x\textasciicircum{}\allowbreak{}2-\allowbreak{}2x-\allowbreak{}1};
in a cyclic ordering the three characters are

\texttt{chi\_\allowbreak{}7(F(Z)\allowbreak{}) = t+\allowbreak{}a+c1*a\allowbreak{}lpha+c\allowbreak{}2*beta\allowbreak{}+c3*ga\allowbreak{}mma},

and its two Galois conjugates. In particular
\texttt{alpha+\allowbreak{}beta+g\allowbreak{}amma=-\allowbreak{}1},
\texttt{alpha\textasciicircum{}\allowbreak{}2+beta\allowbreak{}\textasciicircum{}2+gam\allowbreak{}ma\textasciicircum{}2=5\allowbreak{}}.
The appropriate cyclic ordering is verified against every source product
in the accompanying verifier; permuting the C labels has no effect on
the argument.

We use the standard integral basis
\texttt{O\_\allowbreak{}(Q(alp\allowbreak{}ha))=Z\allowbreak{}[alpha\allowbreak{}]},
with basis
\texttt{1,alph\allowbreak{}a,alph\allowbreak{}a\textasciicircum{}2}. For
example see G. Eric Moorhouse, *Cyclotomic Fields with Applications*,
Theorem 4.8. The accompanying verifier instead tests characteristic
polynomials over Q, so its arithmetic check does not assume an integral
basis theorem.

\hypertarget{only-two-restrictions-can-have-central-dimension-63}{%
\subsection{Only two restrictions can have central dimension
63}\label{only-two-restrictions-can-have-central-dimension-63}}

Lemma. If a central simple has dimension 63, its restriction is one of

\texttt{U = Y+\allowbreak{}B1+B2+\allowbreak{}B3+B4+\allowbreak{}C1+C2+\allowbreak{}C3},

\texttt{V = 1+\allowbreak{}Y+2*C1\allowbreak{}+2*C2+\allowbreak{}2*C3}.

Proof. Item 3 applied to the codegree-9 characters gives

\texttt{t+b-\allowbreak{}a-\allowbreak{}3*b\_\allowbreak{}j == 0\allowbreak{} (mod \allowbreak{}9)}
for every j. (3)

Applied to any codegree-7 character it says
\texttt{(8/\allowbreak{}7)*chi\allowbreak{}\_\allowbreak{}7(F(Z)\allowbreak{})}
is integral. In the integral basis above this forces

\texttt{c1 == \allowbreak{}c2 == \allowbreak{}c3 == \allowbreak{}t+a (m\allowbreak{}od 7)}.
(4)

The dimension identity is

\texttt{t+8a+7\allowbreak{}b+9c=6\allowbreak{}3}, with
\texttt{t in \{\allowbreak{}0,1\}}. (5)

For completeness, here is the entire short reduction of (3)--(5). Write
\texttt{r} for the residue of \texttt{t+a} in
\texttt{\{0,.\allowbreak{}.\allowbreak{}.\allowbreak{},6\}}. Since each
c\_l is nonnegative and congruent to r, the inequality
\texttt{t+8a+2\allowbreak{}7r<=63\allowbreak{}} leaves only the
following aggregate possibilities:

\begin{longtable}[]{@{}rrrrl@{}}
\toprule\noalign{}
t & a & b & c & Consequence \\
\midrule\noalign{}
\endhead
\bottomrule\noalign{}
\endlastfoot
0 & 0 & 9 & 0 & Every b\_j is a multiple of 3 \\
0 & 0 & 0 & 7 & Exactly one c\_l is 7 \\
0 & 1 & 4 & 3 & Every b\_j=1 and every c\_l=1: U \\
0 & 7 & 1 & 0 & (3) requires every b\_j\textgreater=1: impossible \\
1 & 0 & 5 & 3 & (3) requires every b\_j\textgreater=2: impossible \\
1 & 1 & 0 & 6 & Every b\_j=0 and every c\_l=2: V \\
1 & 6 & 2 & 0 & (3) requires every b\_j\textgreater=2: impossible \\
\end{longtable}

In the first row
\texttt{chi\_\allowbreak{}8(F(Z)\allowbreak{})=-\allowbreak{}9}, so its
single contribution to (1) is \texttt{81/\allowbreak{}64>1}. In the
second row the three codegree-7 contributions to (1) sum to
\texttt{alpha\textasciicircum{}\allowbreak{}2+beta\allowbreak{}\textasciicircum{}2+gam\allowbreak{}ma\textasciicircum{}2=5\allowbreak{}>1}.
Both are impossible. Only U and V remain. QED.

\hypertarget{two-dimension-parities-force-prime-power-central-simples}{%
\subsection{Two dimension parities force prime-power central
simples}\label{two-dimension-parities-force-prime-power-central-simples}}

The only possible odd central dimensions are the odd divisors of 504:
\texttt{1,3,7,\allowbreak{}9,21,6\allowbreak{}3}. Dimension 3 is
impossible. Dimension 1 is the unit. At dimensions 7,9,21, the
restrictions are respectively a B\_j, a C\_l, or a sum of three B\_j's;
indeed any occurrence of 1 would put the object in I(1), whose dimension
list has none of 7,9,21. At dimension 63 the preceding lemma applies.

Exactly one central simple has restriction V: it is the unique
63-dimensional summand of I(1). Let m be the number of central simples
with restriction U. Apply (2) to Y and reduce modulo 2. Even-dimensional
central objects contribute zero. Among odd-dimensional central objects,
only the unique V and the m copies of U contain Y, each with
multiplicity one. Therefore

\texttt{0 == 5\allowbreak{}04*8 =\allowbreak{}= 1+m \allowbreak{}(mod 2\allowbreak{})};

hence m is odd.

For each l let n\_l be the number of central simples restricting to
C\_l. Apply (2) to C\_l and again reduce modulo 2. The unique V
contributes \texttt{63*2}, the U columns contribute \texttt{63*m}, and
the dimension-9 columns contribute \texttt{9*n\_\allowbreak{}l}. All
other contributions are even or zero. Thus

\texttt{0 == 5\allowbreak{}04*9 =\allowbreak{}= m+n\_\allowbreak{}l (mod\allowbreak{} 2)}.

Consequently every n\_l is odd. In particular, Z(C) contains a simple of
dimension \texttt{9=3\textasciicircum{}2}.

\hypertarget{contradiction-and-scope}{%
\subsection{Contradiction and scope}\label{contradiction-and-scope}}

The source ring is simple and has no braided categorification. The
latter follows by the standard simple braided dichotomy: a nondegenerate
realization would require
\texttt{7\textasciicircum{}2 | \allowbreak{}504}, while a symmetric
realization would be a finite-group representation ring, which this
exotic ring is not. These ring and braided facts were established in
phase4 and in the published source identifying this ring.

Liu--Palcoux--Ren, arXiv:2102.01663v4, Theorem 5.1, states that the
center of a nonpointed simple complex fusion category is simple exactly
when the category does not admit a braiding. Hence Z(C) is simple. But
Etingof--Nikshych--Ostrik, arXiv:0809.3031, Corollary 7.2, states that
an integral nondegenerate braided category containing a nontrivial
prime-power-dimensional simple contains a nontrivial symmetric
subcategory. Such a subcategory is proper in the nondegenerate Z(C),
contradicting simplicity. This proves the theorem.

The argument applies to arbitrary complex fusion categories because
integrality supplies the positive spherical structure. It never replaces
pseudo-unitarity by unitarity. No associator or indicator classification
is assumed.

\hypertarget{a-forced-lagrangian-subcategory-in-dimension-168}{%
\section{A forced Lagrangian subcategory in dimension
168}\label{a-forced-lagrangian-subcategory-in-dimension-168}}

Theorem. The based ring \texttt{R9\_\allowbreak{}2\_\allowbreak{}1965}
is not the Grothendieck ring of a fusion category over the complex
numbers.

This is a general exclusion, not merely a braided, unitary, solvable,
weakly group-theoretical, or group-theoretical exclusion. The new
argument forces a hypothetical categorification to be group-theoretical
and then applies the independently audited phase-4 exclusion of that
class. In particular, it does not assert that every integral category of
dimension 168 is group-theoretical.

\hypertarget{input-and-categorical-tools}{%
\subsection{Input and categorical
tools}\label{input-and-categorical-tools}}

Use source labels 0 through 8. The positive dimension vector is

\texttt{(1,6,6\allowbreak{},6,6,1\allowbreak{},2,3,3\allowbreak{})},

the dual permutation is
\texttt{(0,2,1\allowbreak{},4,3,5\allowbreak{},6,7,8\allowbreak{})}, and
the global dimension is 168. Put \texttt{g=b5} and \texttt{u=b6}. The
invertible group is \texttt{\{1,g\}=\allowbreak{}C2}. The based adjoint
subring is the entire ring, so the universal grading group is trivial.
These statements, all fusion axioms, and every finite identity below are
checked by
\texttt{verify\allowbreak{}\_\allowbreak{}genera\allowbreak{}l\_\allowbreak{}exclus\allowbreak{}ion.\allowbreak{}py}
against
\texttt{source\allowbreak{}\_\allowbreak{}tensor\allowbreak{}.\allowbreak{}json}.
The tensor's canonical compact-JSON SHA256 is

\texttt{07dac5\allowbreak{}8d848b\allowbreak{}d451d8\allowbreak{}a3d923\allowbreak{}1bfd1c\allowbreak{}23c1b0\allowbreak{}117541\allowbreak{}e6de50\allowbreak{}55fa0a\allowbreak{}b90314\allowbreak{}c4f9}.

Suppose that C categorifies this ring. Since its dimension is integral,
C has its canonical positive spherical structure. Write F: Z(C) → C and
I for its right adjoint. For any irreducible representation E of the
complexified fusion ring, there is a distinguished simple A\_E in I(1)
satisfying

\[
[I(1):A_E]=\dim E,\qquad \dim A_E=\frac{\dim C}{f_E},\qquad\theta_{A_E}=1.
\]

For a scalar ring character χ, the associated central simple satisfies

\[
\frac{S_{A_\chi,Z}}{\dim A_\chi}=\chi(F(Z)). \tag{1}
\]

The induction matrix is a Gram matrix:

\[
B_{ij}=[FI(b_j):b_i]=\sum_Z [F(Z):b_i][F(Z):b_j],
\qquad FI(b_j)=\sum_t b_t b_j b_t^* . \tag{2}
\]

The order of multiplication in (2) matters because the source ring is
noncommutative. Removing any specified central-simple summands from this
Gram decomposition leaves a positive semidefinite matrix.

Primary references for these tools are
\href{https://arxiv.org/abs/1309.4822}{Ostrik, arXiv:1309.4822},
Proposition 2.4, Theorems 2.5 and 2.13, and Example 2.9 with the
discussion preceding Theorem 2.13; also
\href{https://arxiv.org/abs/1611.04486}{Deshpande, arXiv:1611.04486},
Proposition 2.8 and Section 2.7. None requires the fusion ring to be
commutative. Pseudo-unitarity and its canonical structure are
\href{https://math.mit.edu/~etingof/egnobookfinal.pdf}{EGNO, *Tensor
Categories*}, Proposition 9.6.5 and Corollary 9.6.6.

\hypertarget{a-central-tannakian-lift-of-g-is-impossible}{%
\subsection{A central Tannakian lift of g is
impossible}\label{a-central-tannakian-lift-of-g-is-impossible}}

Assume there is a central Tannakian \texttt{Rep(C2\allowbreak{})} whose
nontrivial character forgets to g. The forgetful restriction to this
subcategory is fully faithful. De-equivariantization by its regular
algebra \texttt{1+g} gives a fusion category D of dimension 84 and a
dominant free-module tensor functor Q: C → D (EGNO Theorem 8.23.3).

Tensoring with g has orbits

\texttt{\{0,5\},\allowbreak{} \{1,3\}\allowbreak{}, \{2,4\allowbreak{}\}, \{6\}\allowbreak{}, \{7,8\allowbreak{}\}}.

By free-module adjunction, Q of each simple in a two-element orbit is
simple, with the same dimension as the source. The fixed dimension-two
simple u has \texttt{dim En\allowbreak{}d(Q(u)\allowbreak{})=2}. Hence
it splits into two distinct simples. Their dimensions sum to two and
each is at least one, so both are invertible. Thus D has precisely three
invertibles \texttt{1,a,a²\allowbreak{}}, and its remaining simples are
\texttt{Y,X,Z}, of dimensions \texttt{3,6,6}, respectively. Dominance
ensures that this list is exhaustive. The source images are

\begin{longtable}[]{@{}ll@{}}
\toprule\noalign{}
Source & Q-image \\
\midrule\noalign{}
\endhead
\bottomrule\noalign{}
\endlastfoot
b0,b5 & 1 \\
b1,b3 & X \\
b2,b4 & Z \\
b6 & a+a² \\
b7,b8 & Y \\
\end{longtable}

The order-three group of invertibles fixes Y, X and Z under both left
and right multiplication: it acts on a one-element set of
dimension-three simples and a two-element set of dimension-six simples.
Applying Q to the source products now forces all of the following:

\[
\begin{aligned}
&a^3=1,\quad aY=Ya=Y,\quad aX=Xa=X,\quad aZ=Za=Z,\\
&Y^2=1+a+a^2+2Y,\quad YX=XY=3X,\quad YZ=ZY=3Z,\\
&X^2=2X+4Z,\quad Z^2=4X+2Z,\\
&XZ=ZX=1+a+a^2+3Y+2X+2Z.
\end{aligned} \tag{3}
\]

In particular D is integral. We next give a short theoretical
obstruction to (3).

\hypertarget{the-dimension-84-induction-gram-obstruction}{%
\subsubsection{The dimension-84 induction-Gram
obstruction}\label{the-dimension-84-induction-gram-obstruction}}

In basis \texttt{(1,a,a\allowbreak{}²,Y,X,\allowbreak{}Z)}, (2) gives

\[
B=\begin{pmatrix}
6&3&3&8&4&4\\
3&6&3&8&4&4\\
3&3&6&8&4&4\\
8&8&8&28&12&12\\
4&4&4&12&52&24\\
4&4&4&12&24&52
\end{pmatrix}.
\]

Since (3) is commutative, all six irreducible ring representations are
scalar. Its formal codegrees are
\texttt{3,3,4,\allowbreak{}28,28,\allowbreak{}84}, as verified by the
characteristic polynomial of the Casimir/induction matrix:

\[
\det(fI-B)=(f-3)^2(f-4)(f-28)^2(f-84).
\]

Consequently I\_D(1) has six distinct simple constituents, of dimensions
\texttt{1,3,3,\allowbreak{}21,28,\allowbreak{}28}, each with
multiplicity one. The dimension-one constituent is the central tensor
unit. Let v range over the restriction vectors of the five other
constituents. Each has \texttt{v\_\allowbreak{}0=1}, and

\[
\sum v=(5,3,3,8,4,4). \tag{4}
\]

The two dimension-three constituents each have
\texttt{v\_\allowbreak{}a+v\_\allowbreak{}\{a²\}=2\allowbreak{}}, with
no Y, X or Z coordinates. Thus only two of the total six a/a²
coordinates in (4) remain. For each dimension-28 constituent, its
dimension equation modulo three says
\texttt{v\_\allowbreak{}a+v\_\allowbreak{}\{a²\}≡0\allowbreak{} (mod \allowbreak{}3)}
after subtracting its unit coordinate. The remaining budget is at most
two, so both dimension-28 constituents have zero a and a² coordinates.
The dimension-21 constituent therefore has
\texttt{v\_\allowbreak{}a+v\_\allowbreak{}\{a²\}=2\allowbreak{}}.

Consider the three pairs of a/a² coordinates belonging to dimensions
3,3,21. Each pair sums to two and the three pairs have total
\texttt{(3,3)}. Either all are \texttt{(1,1)}, or they are a permutation
of \texttt{(2,0),\allowbreak{}(1,1),\allowbreak{}(0,2)}. In the latter
case, subtracting their outer products from B gives value
\texttt{6−8=−2\allowbreak{}} on the vector supported as
\texttt{(1,−1)\allowbreak{}} in the a,a² positions. This violates
positive semidefiniteness. Thus all three pairs are \texttt{(1,1)}.

Let t be the Y-coordinate of the dimension-21 constituent, and let y₁,y₂
be those of the two dimension-28 constituents. Equation (4) gives
\texttt{t+y₁+y\allowbreak{}₂=8}. Remove from B the outer products of all
six distinguished restriction vectors, and call the residual matrix M.
Its a,Y principal entries satisfy

\[
M_{aa}=3,\quad M_{aY}=8-t,\quad
M_{YY}=28-t^2-y_1^2-y_2^2\le28-t^2-\frac{(8-t)^2}{2}.
\]

Therefore

\[
\begin{aligned}
\det M_{\{a,Y\}}
&\le84-3t^2-\frac52(8-t)^2\\
&=-\frac{(11t-40)^2}{22}-\frac{36}{11}<0.
\end{aligned}
\]

This contradicts (2). Hence (3) has no fusion categorification, and the
proposed central Tannakian lift of g cannot exist.

For reference only, the full tensor (3) matches the already excluded
lower-census ring \texttt{R6\_\allowbreak{}4\_\allowbreak{}38}, under
permutation \texttt{(0,1,2\allowbreak{},5,3,4\allowbreak{})}. The
preceding proof is independent of its existing exclusion certificate.

\hypertarget{scalar-codegrees-force-a-nontrivial-tannakian-subcategory}{%
\subsection{Scalar codegrees force a nontrivial Tannakian
subcategory}\label{scalar-codegrees-force-a-nontrivial-tannakian-subcategory}}

Set z=−1+√(−7) and z̄=−1−√(−7). Direct substitution in all 81
multiplication equations gives two scalar ring characters

\[
\chi_+=(1,z,\bar z,z,\bar z,1,2,3,3),\qquad
\chi_-=(1,\bar z,z,\bar z,z,1,2,3,3).
\]

Their formal codegrees are both 56. Thus there are two distinct
distinguished central simples A₊,A₋ of dimension three, each occurring
once in I(1), and each having twist one.

Their restrictions contain the unit once. The remaining dimension two
can only be either \texttt{2g} or \texttt{u}: no additional tensor unit
is allowed, because its multiplicity is exactly one. The source
induction matrix has \texttt{B\_\allowbreak{}\{0,5\}=\allowbreak{}1}. A
restriction \texttt{1+2g} would contribute two to this entry and is
impossible. Therefore

\[
F(A_+)=F(A_-)=1+u. \tag{5}
\]

Using (1) and (5), for every σ,τ∈\{+,-\} we obtain

\[
S_{A_\sigma,A_\tau}=3\chi_\sigma(1+u)=9
=\dim A_\sigma\dim A_\tau. \tag{6}
\]

We spell out why (6) gives actual centralization without any assumption
of unitarity. The balancing formula, together with the trivial twists of
Aσ and Aτ, says

\[
9=\sum_W N_{A_\sigma,A_\tau}^{W}\theta_W\dim W,
\qquad
9=\sum_W N_{A_\sigma,A_\tau}^{W}\dim W.
\]

All dimensions are positive, and twists are roots of unity. Equality in
the triangle inequality forces θ\_W=1 for every constituent W. The
balancing identity for the double braiding consequently says that the
double braiding is the identity. Thus A₊ and A₋ centralize one another
and themselves. The fusion subcategory E that they generate is
symmetric. Its twist is identically one: this follows on all tensor
words from balancing and then on their summands and duals. A positive
spherical symmetric fusion category with trivial twist is Tannakian.
Thus

\texttt{E ≃ Re\allowbreak{}p(G)}

for a nontrivial finite group G. Here one can use EGNO's
symmetric-category classification in Section 9.9; the nontrivial super
case has a nontrivial canonical parity twist, which is absent here.

\hypertarget{the-tannakian-subcategory-must-be-lagrangian}{%
\subsection{The Tannakian subcategory must be
Lagrangian}\label{the-tannakian-subcategory-must-be-lagrangian}}

The kernel of \texttt{Inv Z(\allowbreak{}C) → I\allowbreak{}nv C}
consists of half-braidings on the tensor unit, equivalently tensor
automorphisms of the identity functor. These are the characters of the
universal grading group. The source adjoint closure is all of C, so this
kernel is trivial.

It follows that \texttt{Inv E} injects into \texttt{\{1,g\}}. If it were
nontrivial, E would contain a Tannakian \texttt{Rep(C2\allowbreak{})}
whose nontrivial character forgets to g. Section 2 excluded precisely
this situation. Therefore \texttt{Inv E} is trivial: G has no nontrivial
one-dimensional representations and is perfect.

Let E′ be its centralizer in Z(C). The centralizer dimension theorem and
de-equivariantization give a fusion category
\texttt{(E′)\_\allowbreak{}G} of dimension

\[
\operatorname{FPdim}((E')_G)=\frac{168^2}{|G|^2}.
\]

This rational number is an algebraic integer, so it is an integer. A
rational number whose square is an integer is itself an integer; hence
\texttt{|G|} divides 168.

Every group of proper-divisor order of 168 is solvable. Indeed, the
orders involving at most two primes are covered by Burnside's theorem.
The only proper divisors involving all three primes 2,3,7 are 42 and 84;
in either order the Sylow-seven number divides 6 or 12 and is congruent
to one modulo seven, so it is one. The quotient by this normal C7 has
order 6 or 12 and is solvable. Since a nontrivial perfect group is not
solvable, necessarily \texttt{|G|=16\allowbreak{}8}.

Therefore \texttt{FPdim \allowbreak{}E=168} and
\texttt{FPdim \allowbreak{}Z(C)=1\allowbreak{}68²}. Since E is
symmetric, the centralizer dimension formula now yields \texttt{E=E′}:
it is a Lagrangian subcategory. EGNO Corollary 9.13.6 implies that C is
group-theoretical.

\hypertarget{contradiction-with-the-complete-group-theoretical-obstruction}{%
\subsection{Contradiction with the complete group-theoretical
obstruction}\label{contradiction-with-the-complete-group-theoretical-obstruction}}

Phase 4 already excluded every group-theoretical categorification of
\texttt{R9\_\allowbreak{}2\_\allowbreak{}1965}. The proof is Section 5
of
\texttt{phase4\allowbreak{}/\allowbreak{}integr\allowbreak{}al/\allowbreak{}GROUP\_\allowbreak{}THEORE\allowbreak{}TICAL\_\allowbreak{}OBSTRU\allowbreak{}CTIONS\allowbreak{}.\allowbreak{}md},
with independent replay in
\texttt{phase4\allowbreak{}/\allowbreak{}audit\_\allowbreak{}integr\allowbreak{}al\_\allowbreak{}gt/\allowbreak{}AUDIT.\allowbreak{}md}.

Briefly, the all-twists projective-character reduction exhausts all 57
groups of order 168 and all 1,751 subgroup conjugacy classes. Only
\texttt{(PSL(2\allowbreak{},7),S4\allowbreak{})}, for the two S4
conjugacy classes, survives. Restriction on third cohomology is
injective: its kernel is seven-primary, whereas the Sylow-seven
normalizer acts without fixed points on \texttt{H⁴(C7,\allowbreak{}Z)}.
Thus every admissible associator is globally trivializable. Both classes
in \texttt{H²(S4,\allowbreak{}C*)} extend from
\texttt{H²(PSL\allowbreak{}(2,7),\allowbreak{}C*)} and can be removed by
a global tensor autoequivalence. The resulting untwisted tensors are
\texttt{R9\_\allowbreak{}2\_\allowbreak{}1557}, not
\texttt{R9\_\allowbreak{}2\_\allowbreak{}1965}. Their duality types
already differ (seven versus five self-dual simples).

The group-theoretical exclusion is incompatible with Section 4. This
proves the theorem.

\hypertarget{an-isotropic-orbifold-model-for-r921848}{%
\section{\texorpdfstring{An isotropic orbifold model for
R9\emph{2}1848}{An isotropic orbifold model for R921848}}\label{an-isotropic-orbifold-model-for-r921848}}

R9\_2\_1848 is modelled. R9\_2\_1975 remains unresolved.

The new model is an application of an existing, general orbifold theorem
to an existing unitary near-group category. It does not require a new
associator ansatz or the unproved assertion that every spherical
category is unitary.

\hypertarget{construction-and-complete-fusion-rules}{%
\subsection{Construction and complete fusion
rules}\label{construction-and-complete-fusion-rules}}

Let N be Budinski's near-group category J\_18\^{}1 of type C18+18. Its
simple objects consist of the group G=C18 and one object rho with

\begin{verbatim}
rho\^2 = sum\_(g in C18) g + 18 rho,
FPdim(rho) = 9 + 3 sqrt(11).
\end{verbatim}

The source exists as a unitary category by the Cuntz algebra
construction described in Budinski, Theorem 3.1.1 and Table 3.1 (the
four rows for n=18). Table 3.1 displays rounded phase coordinates, but
section 3.3.2 explicitly explains the exact algebraic certification of
these solutions. We use the source's existence result as a literature
dependency. This continuation does not claim to independently
reconstruct or replay its complete associator.

For J\_18\^{}1 the nondegenerate symmetric pairing and quadratic
refinement are

\begin{verbatim}
<x,y> = exp(2 pi i xy/18),
a(x) = exp(-pi i x\^2/18).
\end{verbatim}

Take H=\{0,6,12\}. Then

\begin{verbatim}
H\^perp = {0,3,6,9,12,15},
H subset H\^perp,  a|H=1,
G/H = C6,  H\^perp/H = C2,  G/H\^perp = C3.
\end{verbatim}

Apply Izumi, Theorem 12.5, together with the irreducible decomposition
and multiplication formulas immediately preceding it in section 12.1.
This theorem covers isotropic H; it does not require H to be Lagrangian.
The resulting unitary orbifold has six distinct invertibles indexed by
G/H, and three distinct noninvertibles indexed by G/H\^{}perp. With g a
generator of C6 and X the distinguished component, its complete
presentation is

\begin{verbatim}
g\^6=1,  g\^3 X=X,  Xg=g\^-1 X,
X\^2=1+g\^3+2(X+gX+g\^2X).
\end{verbatim}

Here the theorem's shift g\_a may be chosen to be zero because
a\textbar H=1. Its noninvertible coefficient is
\textbar H\^{}perp/H\textbar=2. In particular, FPdim(X)=3+sqrt(11), all
three noninvertibles are self-dual, and the category has rank nine and
maximum fusion multiplicity two. The construction supplies associativity
at the categorical level; checking the based tensor alone would not
supply that missing premise.

The construction is a special case of the unified formula

\begin{verbatim}
sigma\^2 = sum\_(k in H\^perp/H) alpha\_(k-g\_a)
          + |H\^perp/H| sum\_(q in G/H\^perp) alpha\_q sigma,
alpha\_g sigma = sigma alpha\_-g.
\end{verbatim}

Thus the model uses a known theoretical construction and a known source
category, rather than claiming a new family of associators.

\hypertarget{exact-census-binding}{%
\subsection{Exact census binding}\label{exact-census-binding}}

Order the model basis as

\begin{verbatim}
[1,g,g\^2,g\^3,g\^4,g\^5,X,gX,g\^2X].
\end{verbatim}

The map into the census tensor basis is

\begin{verbatim}
[0,1,4,8,3,2,5,7,6].
\end{verbatim}

All 729 coefficients agree with R9\_2\_1848. The source tensor SHA256 is
\texttt{426814\allowbreak{}07529a\allowbreak{}88aaaf\allowbreak{}34a909\allowbreak{}975776\allowbreak{}d8d100\allowbreak{}8ee8a9\allowbreak{}9b6e23\allowbreak{}1263dc\allowbreak{}b1aae9\allowbreak{}ff5c}.
The self-contained record
\texttt{1848\_\allowbreak{}model.\allowbreak{}json} includes both
tensors and the permutation. Run

\begin{verbatim}
python verify\_1848\_model.py
\end{verbatim}

to verify the transport, the subgroup and root-of-unity exponent
conditions, associativity, duality, and the dimension identities using
exact integers. The receipt is
\texttt{1848\_\allowbreak{}model\_\allowbreak{}verifi\allowbreak{}cation\allowbreak{}.\allowbreak{}json}.
Its scope explicitly separates the new exact transport from the cited
near-group existence theorem.

\hypertarget{primary-sources-and-scope}{%
\subsection{Primary sources and scope}\label{primary-sources-and-scope}}

\begin{enumerate}
\def\labelenumi{\arabic{enumi}.}
\tightlist
\item
  P. G. Budinski, *Exotic Fusion Categories and Their Modular Data*,
  master's thesis, University of Alberta (2021), Theorem 3.1.1, Table
  3.1 (n=18), section 3.3.2 (exactness, printed pp.46--47). DOI:
  \url{https://doi.org/10.7939/r3-7p7k-fg18}. Institutional PDF:
  \url{https://ualberta.scholaris.ca/bitstreams/3f73f320-a2a2-4184-ad81-d8708764f725/download}.
\item
  M. Izumi, *A Cuntz algebra approach to the classification of
  near-group categories*, arXiv:1512.04288v2, section 12.1, Theorem 12.5
  and the preceding action formulas (printed pp.107--111):
  \url{https://arxiv.org/abs/1512.04288v2}.
\end{enumerate}

Source URLs and hashes are in
\texttt{source\allowbreak{}s.\allowbreak{}json}. Downloaded primary
papers and extracted text are held in the local research cache; they are
not included in the verified baseline archive or the final research
package. In particular, Izumi's cached PDF and text under
\texttt{source\allowbreak{}s} are URL/hash dependencies, not bundled
replay inputs. The exact tensor replays do not read these cached papers.
Later citations to the thesis in Yu's 2024 paper and Yu--Zhang's notes
support that its solutions are used as existence results, but the proof
dependency above is the primary thesis itself.

Izumi works in the unitary C* setting. That is sufficient here because
the conclusion is the existence of one unitary model. No classification
theorem for arbitrary original spherical categories is inferred.

}

\section{What remains unresolved}
\begin{center}\begin{tabular}{ll}\toprule Type & Suffixes of unresolved $R9\_2\_n$\\\midrule
Integral & 12, 941, 959, 1965\\
Strictly weakly integral & 585, 746, 1025, 1209, 1254,\\
 & 1303, 1709, 1759, 1988\\
Other quadratic & 137, 1607, 1848, 1975, 1976\\
\bottomrule\end{tabular}\end{center}

The four remaining rings are $R9\_2\_137,1607,1975,1976$, all with
dimension type $(1^4,2,(4+3\sqrt2)^4)$ and global dimension
$144+96\sqrt2$. General, pivotal, spherical, pseudo-unitary and unitary
existence remain undecided. The audited sphericalization and Galois gates
equate general existence with pseudo-unitary existence; they do not make
every original category pseudo-unitary or unitary. Braided categorification
is excluded for all four.

The stronger center, TY-sign and Clifford restrictions are developed in
Section~\ref{sec:quadratic-frontier-retry}. There are three necessary
identity-action D8 types in the union: zero or reflection shift for the
commutative sources, and reflection or rotation for the noncommutative
source. Only the reflection type overlaps. Their based tensors do not
supply compatible associators. The realized auxiliary orbifolds still
leave the source categorifications open.

\section{Unified restrictions on the four quadratic cases}
\label{sec:quadratic-frontier-retry}
Throughout this section $d=4+3\sqrt2$ and $D=144+96\sqrt2$.
The four quadratic rings have dimension type $(1^4,2,d^4)$, and
$N_{\mathbb Q(\sqrt2)/\mathbb Q}(D)=2304=2^8 3^2$.
All spherical and positive-center arguments use the already audited
sphericalization and Galois-existence reduction; no claim is made that
every original category is unitary or pseudo-unitary.

\subsection{The integral part of the center is pointed}
\begin{proposition}
For a complex categorification of any of the four quadratic rings,
the center contains no nontrivial Tannakian subcategory, and every
integral-dimensional central simple is invertible. The integral center
is Vec or sVec for the three commutative rings, and Vec for $R9\_2\_1975$.
\end{proposition}
\begin{proof}
The established TY module reduction gives the nonalternating bicharacter
on $V_4$. In a commutative source, let $h,h'$ stabilize its two distinct
large-object orbits. Their product $g$ is isotropic and acts freely.
The adjoint subcategory is the full source, so the universal grading
is trivial and forgetting central invertibles is injective.
The TY central-invertible formulas give self-braiding $+1$ for a lift
of $g$, and $-1$ for a lift of $h$ or $h'$. Thus a lift of $g$ would
supply a central Rep$(C_2)$.

Central de-equivariantization would have four invertibles and two
self-dual objects $X,Y$ of dimension $d$. Its rules are
\begin{gather*}
 hX=X,\quad hY=Y,\quad uX=Y,\quad uY=X,\\
 X^2=Y^2=1+h+4X+4Y,\qquad XY=u+hu+4X+4Y.
\end{gather*}
The alternatives $u^2=1$ and $u^2=h$ are exactly the already excluded
project tensors $R6\_4\_174$ and $R6\_4\_6$. All $216$ coefficients
of each identification agree. Their general exclusions are inherited
lower-rank theorems, not deductions from failure to find a model.
Hence $g$ has no central lift. There is at most one fermionic lift of
$h,h'$, since two would lift $g$. This proves the pointed-center bound.

For the noncommutative source, the corresponding hypothetical central
quotients are the excluded $R6\_4\_1$ or $R6\_4\_155$ with a pointed
permutation that cannot lift to a tensor autoequivalence. The latter
obstruction follows from the distinct uniquely lifted pointed twists
of $R6\_4\_155$. The exact source-to-quotient maps and both alternatives
are checked in the proof record. The pointed center is therefore Vec.

If Rep$(H)$ were Tannakian in the center, FP-dimension divisibility
and the norm of $D$ would force the primes of $|H|$ to lie in $\{2,3\}$.
Burnside makes $H$ solvable. A nontrivial solvable group has a nonunit
one-dimensional representation, which would give a central boson,
contradicting the preceding bounds.

Each central dimension is $a+bd$ with nonnegative integral $a,b$.
Thus its integral-dimensional simples are exactly those forgetting into
the integral TY subcategory. They form an integral braided fusion
subcategory whose dimension has prime support in $\{2,3\}$. The
categorical Burnside theorem makes it solvable, hence weakly
group-theoretical. With no Tannakian subcategory it is weakly
anisotropic, so Natale's integral weakly-anisotropic theorem makes it
pointed. This applies to the integral subcategory, not to the ambient
nonintegral center.
\end{proof}

The categorical inputs are the TY formulas of Gelaki--Naidu--Nikshych,
Proposition 4.1~\cite{RetryGNN}; ENO FP-dimension divisibility; Etingof--Nikshych--Ostrik,
\emph{Weakly group-theoretical and solvable fusion categories},
Theorem 1.6~\cite{RetryENO}; and Natale, \emph{The core of a weakly group-theoretical
braided fusion category}, Theorem 4.1~\cite{RetryNatale}. Their hypotheses and the exact
lower-rank dependencies are recorded in the linked proof package.
In particular no Morita representative has a nontrivial faithful
finite-group grading, since such a grading would give central Rep$(H)$.

\subsection{TY saturation and its sign}
Let $Y$ be the dimension-two TY object and $\epsilon$ the sign of its
parameter in a positive spherical representative. Shimizu's formula~\cite{RetryShimizu} gives
\[
 (\nu_1(Y),\ldots,\nu_8(Y))=(0,\epsilon,0,0,0,\epsilon,0,2).
\]
The center induction formula, $\nu_8(Y)=\FPdim(Y)$ and positivity force
all constituents of $I(Y)$ to have eighth-root twists. Fourier inversion
then fixes the weighted fibers. Each primitive-eighth fiber has weight
$D/4$; its only possible dimension pairs are
\[
 (2+d,2+7d)\quad\hbox{or}\quad(2+4d,2+4d).
\]
Combining this complete support condition with the integral-center
obstruction and character/induction equations gives exact separating
certificates for $\epsilon=+1$ in $R1607,R1975,R1976$. For R137 these support conditions alone do not decide the sign; Theorem~\ref{thm:r137-positive-ty-exclusion} excludes its positive sign using additional integral Verlinde restrictions.
An exact witness is evidence about those systems, not a category.

\subsection{A general sign-twisted module-functor criterion}
\begin{proposition}
Let $\mathcal D$ be a fusion subcategory, $\mathcal M$ an indecomposable
left $\mathcal D$-component in a fusion category, and $X$ a simple in it.
Suppose $h$ is an order-two invertible in $\mathcal D$,
$\underline{\operatorname{End}}_{\mathcal D}(X)=1+h$, $Xh\cong X$, and
right tensoring by $h$ moves another simple of $\mathcal M$.
Then the normalized left and right stabilizer maps on $X$ anticommute.
\end{proposition}
\begin{proof}
The actual internal end $B=1+h$ has its subgroup-algebra structure and
identifies $\mathcal M$ with right $B$-modules in $\mathcal D$, taking
$X$ to $B$. The functor $(-)h$ is a left $\mathcal D$-module functor.
The module-functor correspondence~\cite{RetryEGNO} represents it by a $B$-bimodule $P$. Since $Xh\cong X$,
its underlying right $B$-module is $B$. Its left action is an algebra
automorphism of $B$: the unit is fixed and the $h$-summand maps to
$h$ or $-h$. These exhaust the possibilities because
$\operatorname{Hom}_B(hB,B)$ is one-dimensional and $h^2=1$.
The positive action gives the regular bimodule and hence the identity
module functor, impossible because another simple is moved. The action
is therefore the negative one. Transporting it through $r:Xh\to X$
gives, with $l:hX\to X$,
\[
 r(l\otimes1_h)=-l(1_h\otimes r),
\]
with ambient associators understood. This is the required categorical
anticommutator. It uses no central lift of the based-central $h$.
\end{proof}

For each of the three commutative sources, the two pointed orbits have
different stabilizers, so right $h$ fixes the first and swaps the second.
The criterion applies to every large simple. Exact source tensors give
$\dim\operatorname{Hom}(1,X^3)=2$ in all twelve applications. The
Clifford rotation theorem then forces
\begin{equation}\label{eq:source-cubic-clifford}
 \nu_3(X)\in\{-1,-\zeta_3,-\zeta_3^2\}.
\end{equation}
The normalized relative trace to $\mathbb Q(\sqrt2)$ is consequently
$-1$ or $1/2$. No reality assertion is inferred from self-duality.
The noncommutative source has different left and right stabilizers;
this source-local criterion is not applied to it.

\subsection{The all-cochain Morita restriction}
The rank-twelve Morita reductions have pointed group $D_8$ and four
objects of dimension $d$. There are three types in their union:
commutative sources allow duality shift zero or a reflection coset;
$R1975$ allows reflection or rotation. Only the reflection type overlaps.
The earlier nonidentity quotient action and extreme $R1975$ branch
are excluded by their separate associator and pointed-automorphism
obstructions.

The forced cyclic restrictions determine the pointed $D_8$ associator
class: it restricts to class two on the rotation $C_4$, trivially on
the even $V_4$, and to the TY sign on the odd reflection groups.
The cohomology calculation of Goff--Mason--Ng~\cite[Lemmas 4.4--4.5]{RetryGMN}
shows that these cyclic restrictions detect the full class.

For a large simple, the pointed-bimodule stabilizer is
$L=D_8\times_{V_4}D_8\cong D_8\times C_2$. An explicit stabilizer
cochain, checked against every coboundary equation, has mixed pairing
$-1$ on its pure left and right central involutions $h_L,h_R$.
This is independent of the chosen solution. Indeed any other solution
differs by a $2$-cocycle $\gamma$. Its commutator with the central
$h_R$ is a character of $L$, hence is trivial on $[L,L]$; but
$h_Lh_R\in[L,L]$ and its pairing with $h_R$ equals that of $h_L$.
Thus changing the cochain cannot change the sign.

The nondegenerate stabilizer pairing makes this four-object pointed
bimodule invertible~\cite[Proposition 2.7]{RetryVV}. The cubic Clifford theorem gives
\eqref{eq:source-cubic-clifford} for each of its four large objects.
For pseudo-unitary Morita partners the weighted total indicator
\[
 v_3(\mathcal C)=\sum_T\FPdim(T)\nu_3(T)
 =\frac1D\sum_{Z\in\operatorname{Irr}\mathcal Z(\mathcal C)}
                 \FPdim(Z)^2\theta_Z^3
\]
is Morita invariant by the center indicator formula~\cite{RetryNS}. The nonunit invertibles and TY object have third
indicator zero. Subtracting the common unit and dividing by $d$
therefore transfers the sum of four large-object indicators. Its
projected value belongs to
\[
 \{-4,-5/2,-1,1/2,2\}.
\]
This covers $R1975$ as well. Neither bimodule invertibility nor
based centrality supplies a central half-braiding on the full category.

\subsection{Cubic branch restrictions and exponent consequence}
The cubic tests alone have the following feasible projected choices in their necessary center systems; these are intermediate relaxations, not the final status of a TY sign:
\begin{center}\begin{tabular}{lcl}
\toprule Source & TY sign & Four projected third indicators\\\midrule
$R137$ & $-$ & $(-1,-1,-1,-1)$\\
$R137$ & $+$ & $(1/2,1/2,1/2,1/2)$\\
$R1607$ & $-$ & $(-1,-1,-1,-1)$\\
$R1976$ & $-$ & $(1/2,1/2,1/2,1/2)$\\\bottomrule
\end{tabular}\end{center}
The $R1975$ height-six system with the aggregate condition also has
an exact witness at projected sum $2$.

For positive-sign $R137$ and negative-sign $R1976$, all other discrete
assignments are excluded: $25$ exact separating vectors and five
verified source relabellings cover all $15$ alternatives in each case.
The relabellings fix the distinguished pointed and TY objects and are
used to relabel a hypothetical categorification; no tensor
autoequivalence lifting them is assumed. Both saved matrices were
independently reconstructed entrywise.

In these two conditional branches, every actual third indicator is
$-\zeta_3$ or $-\zeta_3^2$, and $9$ divides the FS exponent. Otherwise
cubing every central twist would remove its $3$-primary part. Since
all central dimensions lie in $\mathbb Q(\sqrt2)$, the induction formula
would put $\nu_3(X)$ in a cyclotomic field of conductor coprime to $3$,
which cannot contain a nonrational cube root of unity. This contradicts
the forced indicators.

The complete proofs, explicit cochains, bounded-system semantics and
portable exact replays are curated in
\path{evidence/unresolved_retry_20261024/README.md}.
These cubic restrictions supply necessary premises, not associators or categorifiability. The following section excludes three source rings, and Theorem~\ref{thm:r137-positive-ty-exclusion} excludes the positive TY sign of the remaining ring.



\section{Exact evidence and reproduction}\label{sec:replay}
The public source is
\url{https://github.com/sebastienpalcoux/grothendieck-ring-classification/tree/main/evidence/rank9_m2_20261024}.
It exposes the accepted ledger, positive models, exact open tensors and
readable proofs. Its checksummed replay archive retains every source
cochain, certificate, constructor, independent audit and historical source
context. Only \path{phase5/results/} is the current mathematical partition.
Historical ledgers remain explicitly superseded; their unsupported half-$E_6$
claims are not used by the corrected verifier.

After unpacking the archive and installing its requirements, run
\begin{quote}\ttfamily
python3 verify.py\\
python3 verify.py --frontier
\end{quote}
The first command verifies archive hashes, accepted finite evidence, every
input identity and all 36 exact model transports. The second adds the
auxiliary Morita and algebraic-orbifold checks. The earlier research
snapshots and their supersession records remain preserved. Each categorical
argument has a separate review; computation checks its finite premises,
not the cited categorical theorems. Independent human verification and
formalization remain pending.

The fusion census repository publishes the same 2,225 supplied tensors
through multiplicity two and the corresponding producer-run qualification.
The public Atlas uses one rank/multiplicity table and one case directory
for all ranks. Every M2 case has its source identity, current decision,
proof links and a downloadable Python tensor. The four open cases
retain their separate categorical-scope statements.

\begin{thebibliography}{99}
\bibitem{ENO} P.~Etingof, D.~Nikshych, V.~Ostrik, \emph{On fusion categories}, Ann. Math. 162 (2005), 581--642; \href{https://arxiv.org/abs/math/0203060}{arXiv:math/0203060}.
\bibitem{Ostrik} V.~Ostrik, \emph{Pivotal fusion categories of rank 3}, Mosc. Math. J. 15 (2015), 373--396; \href{https://arxiv.org/abs/1309.4822}{arXiv:1309.4822}.
\bibitem{Deshpande} T.~Deshpande, \emph{On centers of bimodule categories and induction-restriction functors}, \href{https://arxiv.org/abs/1611.04486}{arXiv:1611.04486}, Proposition 2.8 and Section 2.7.
\bibitem{DMNO} A.~Davydov, M.~M\"uger, D.~Nikshych, V.~Ostrik, \emph{The Witt group of non-degenerate braided fusion categories}, J. Reine Angew. Math. 677 (2013), 135--177; \href{https://arxiv.org/abs/1009.2117}{arXiv:1009.2117}.
\bibitem{HH} T.~Hagge, S.-M.~Hong, \emph{Some non-braided fusion categories of rank three}, Commun. Contemp. Math. 11 (2009), 615--637; \href{https://arxiv.org/abs/0704.0208}{arXiv:0704.0208}.
\bibitem{ALW} D.~Aasen, E.~Lake, K.~Walker, \emph{Fermion condensation and super pivotal categories}, J. Math. Phys. 60 (2019), 121901; \href{https://arxiv.org/abs/1709.01941}{arXiv:1709.01941}, Section 7.
\bibitem{ChengNg} Y.-L.~Cheng, S.-H.~Ng, \emph{On Hopf algebras of dimension $4p$}, J. Algebra 328 (2011), 399--419; \href{https://arxiv.org/abs/1005.1267}{arXiv:1005.1267}.
\bibitem{Natale} S.~Natale, \emph{On the semisolvability of semisimple Hopf algebras of low dimension}, \href{https://arxiv.org/abs/math/0305038}{arXiv:math/0305038}, Section 5.2.
\bibitem{AN} N.~Andruskiewitsch, S.~Natale, \emph{Examples of self-dual Hopf algebras}, J. Math. Sci. Univ. Tokyo 6 (1999), 181--215; \url{https://www.ms.u-tokyo.ac.jp/journal/pdf/jms060108.pdf}.
\bibitem{DR} A.~Davydov, I.~Runkel, \emph{$\Z/2\Z$-extensions of Hopf algebra module categories by their base categories}, \url{https://www.math.uni-hamburg.de/home/runkel/PDF/hext.pdf}, Corollary 3.19 and Section 3.6.
\bibitem{GHR} C.~Galindo, S.-M.~Hong, E.~C.~Rowell, \emph{Generalized and quasi-localizations of braid group representations}, \href{https://arxiv.org/abs/1105.5048}{arXiv:1105.5048}, Theorem 5.20.
\bibitem{ABDP} M.~Alekseyev, W.~Bruns, J.~Dong, S.~Palcoux, \emph{Classifying integral Grothendieck rings up to rank 5 and beyond}, \href{https://arxiv.org/abs/2507.07023}{arXiv:2507.07023}, Question 1.14.
\bibitem{RetryGNN} S.~Gelaki, D.~Naidu, D.~Nikshych, \emph{Centers of graded fusion categories}, \href{https://arxiv.org/abs/0905.3117}{arXiv:0905.3117}, Proposition 4.1.
\bibitem{RetryENO} P.~Etingof, D.~Nikshych, V.~Ostrik, \emph{Weakly group-theoretical and solvable fusion categories}, Adv. Math. 226 (2011), 176--205; \href{https://arxiv.org/abs/0809.3031}{arXiv:0809.3031}, Theorem 1.6.
\bibitem{RetryNatale} S.~Natale, \emph{The core of a weakly group-theoretical braided fusion category}, \href{https://arxiv.org/abs/1704.03523}{arXiv:1704.03523}, Theorem 4.1.
\bibitem{RetryShimizu} K.~Shimizu, \emph{Frobenius--Schur indicators in Tambara--Yamagami categories}, \href{https://arxiv.org/abs/1005.4500}{arXiv:1005.4500}, Theorem 3.4.
\bibitem{RetryGMN} C.~Goff, G.~Mason, S.-H.~Ng, \emph{On the gauge equivalence of twisted quantum doubles of elementary abelian and extra-special 2-groups}, \href{https://arxiv.org/abs/math/0603191}{arXiv:math/0603191}, Lemmas 4.4--4.5.
\bibitem{RetryVV} L.~Vainerman, J.-M.~Vallin, \emph{On $\mathbb Z/2\mathbb Z$-extensions of pointed fusion categories}, \href{https://arxiv.org/abs/1210.1732}{arXiv:1210.1732}, Proposition 2.7.
\bibitem{RetryNS} S.-H.~Ng, P.~Schauenburg, \emph{Frobenius--Schur indicators and exponents of spherical categories}, \href{https://arxiv.org/abs/math/0601012}{arXiv:math/0601012}, Theorem 4.1.
\bibitem{RetryEGNO} P.~Etingof, S.~Gelaki, D.~Nikshych, V.~Ostrik, \emph{Tensor Categories}, AMS Mathematical Surveys and Monographs 205 (2015).
\end{thebibliography}
\end{document}
